\documentclass[11pt,a4paper,openany]{report}
\usepackage[T1]{fontenc}
\usepackage[utf8]{inputenc}
\usepackage{lmodern}
\usepackage[margin=23mm]{geometry}
\usepackage{amsmath,amssymb,amsthm,mathtools}
\usepackage{longtable,booktabs,array,calc}
\newcounter{none}
\usepackage{graphicx,xcolor,tikz}
\usepackage{microtype}
\usepackage[unicode,breaklinks,colorlinks,linkcolor=blue!45!black,urlcolor=blue!45!black]{hyperref}
\urlstyle{same}
\providecommand{\tightlist}{\setlength{\itemsep}{0pt}\setlength{\parskip}{0pt}}
\providecommand{\passthrough}[1]{#1}
\providecommand{\NL}{\operatorname{NL}}
\providecommand{\colim}{\operatorname*{colim}}
\setlength{\parindent}{0pt}
\setlength{\parskip}{.4em}
\setlength{\emergencystretch}{3em}
\setcounter{tocdepth}{1}
\setcounter{secnumdepth}{0}
\renewcommand{\chaptername}{Lesson}
\DeclareUnicodeCharacter{03BB}{\ensuremath{\lambda}}
\DeclareUnicodeCharacter{0393}{\ensuremath{\Gamma}}
\DeclareUnicodeCharacter{2192}{\ensuremath{\rightarrow}}
\DeclareUnicodeCharacter{2011}{-}
\DeclareUnicodeCharacter{202F}{\,}
\DeclareUnicodeCharacter{2212}{\ensuremath{-}}
\DeclareUnicodeCharacter{03B1}{\ensuremath{\alpha}}
\DeclareUnicodeCharacter{03B2}{\ensuremath{\beta}}
\DeclareUnicodeCharacter{03BC}{\ensuremath{\mu}}
\DeclareUnicodeCharacter{03A6}{\ensuremath{\Phi}}
\DeclareUnicodeCharacter{2260}{\ensuremath{\ne}}
\DeclareUnicodeCharacter{00D7}{\ensuremath{\times}}
\DeclareUnicodeCharacter{2082}{\ensuremath{_2}}
\DeclareUnicodeCharacter{2228}{\ensuremath{\vee}}
\DeclareUnicodeCharacter{2032}{\ensuremath{^{\prime}}}
\DeclareUnicodeCharacter{220E}{\ensuremath{\square}}
\title{Roots and reductive groups of rank one}
\author{GPT-6.1 Sol (OpenAI)\\Codex, Ultra setting}
\date{October 2026}
\begin{document}
\hypersetup{pageanchor=false}
\maketitle
\pagenumbering{roman}
\hypersetup{pageanchor=true}
\noindent Written and self-checked by the writing AI. No human or independent review is claimed.
Original contributions are dedicated under CC0 1.0. Cited works retain their own licences.
\par\medskip
\noindent Online course: \url{https://kokunoyumeto.github.io/open-math-courses/courses/AG-RG/}.
\begingroup\small\tableofcontents\endgroup
\clearpage
\pagenumbering{arabic}
\setcounter{chapter}{2}
\chapter{Roots and reductive groups of rank one}\label{AG-RG-03-roots-and-reductive-groups-of-rank-one}

\emph{Written by GPT-6.1 Sol (OpenAI) in Codex, Ultra setting, October 2026. Self-checked by the writing AI, GPT-6.1 Sol (OpenAI). Original text: public domain (CC0). The normalization lemma in Section 5 follows the Stacks Project's proof, cited at the end.}

The diagonal entries of a matrix act on its off-diagonal entries through ratios. Those ratios are the roots of a general linear group. For an arbitrary reductive group, roots are still characters of a maximal torus, but their one-dimensional spaces need not be globally trivial over the base. This lesson constructs the corresponding groups without using a power-series exponential, and shows why every root carries a copy of the rank-one geometry of \(\operatorname{SL}_2\).

We use torus conjugacy, torus deformation, and the centralizer theorem from the preceding two lessons. A reductive \(S\)-group is smooth affine of finite presentation with connected reductive geometric fibres. A semisimple \(S\)-group has semisimple geometric fibres. Statements about line bundles and subgroup schemes are made over arbitrary schemes, including nonreduced ones.

\section{1. The centre and the two ranks}\label{AG-RG-03-the-centre-and-the-two-ranks}

Let \(T\) be a split maximal torus of \(G\), and, on an open set where the ranks are constant, write

\[
\mathfrak g=\mathfrak t\oplus\bigoplus_{\alpha\in\Phi}\mathfrak g_\alpha.
\]

For now \(\Phi\) means the nonzero weights; we will prove that they form a reduced root system. Put

\[
Z=\bigcap_{\alpha\in\Phi}\ker(\alpha:T\to\mathbf G_m).
\]

The group \(Z\) is of multiplicative type. It is precisely the scheme-theoretic centre of \(G\). Indeed, its conjugation action on \(\mathfrak g\) is trivial. Exactness of invariants for a multiplicative-type group makes \(C_G(Z)\) smooth with Lie algebra \(\mathfrak g\). Its geometric fibres are therefore all of \(G\); the smooth fibrewise isomorphism argument of the preceding lesson gives \(C_G(Z)=G\). Conversely any central point belongs to \(C_G(T)=T\) and acts trivially on every weight summand, hence lies in \(Z\). These constructions agree on overlaps and descend when \(T\) is only étale locally split.

The centre need not be smooth: \(Z(\operatorname{SL}_n)=\mu_n\), including when the characteristic divides \(n\). Its largest subtorus is called the \textbf{radical}, denoted \(\operatorname{Rad}(G)\). For a group of multiplicative type with character sheaf \(M\), this subtorus has character sheaf \(M/M_{\mathrm{tors}}\). Thus its rank is locally constant and its formation commutes with base change. The \textbf{rank} of \(G\) is the dimension of a maximal torus; its \textbf{semisimple rank} is

\[
\operatorname{rank}(G)-\dim\operatorname{Rad}(G).
\]

Over an algebraically closed field, the radical also equals the largest smooth connected solvable normal subgroup. Its unipotent radical is normal in \(G\) and is trivial by reductivity, so it is a torus; rigidity of torus automorphisms makes it central. Conversely a central torus is connected solvable normal. The quotient by this torus is semisimple: the inverse image of a connected solvable normal subgroup of the quotient would enlarge the radical.

The derived group will be constructed for arbitrary bases after the pinned existence theorem in \href{https://kokunoyumeto.github.io/open-math-courses/courses/AG-RG/AG-RG-05.html\#the-derived-group}{Pinnings and the classification of split reductive groups}. It is semisimple and the multiplication map

\[
\operatorname{Rad}(G)\times G_{\mathrm{der}}\longrightarrow G
\]

is a central isogeny. Its construction uses the root groups established here. We do not assume a relative derived group in constructing these groups. Central quotients are constructed by graded invariant algebras in Theorem 1.1 below, without using coroots. In rank one we will construct the simply connected cover explicitly below.

\subsection{Central quotients before the rank-one argument}\label{AG-RG-03-central-quotients-before-the-rank-one-argument}

The rank-one argument below needs the affine quotient even when the centre is not smooth. The invariant-algebra construction is independent of coroots and root groups.

\textbf{Theorem 1.1 (affine central quotient).} Let \(G\) be a reductive group over a scheme \(S\), and let \(H\subset Z(G)\) be a closed subgroup of multiplicative type of finite presentation. The fppf quotient \(G/H\) is an affine reductive \(S\)-group. Formation of the quotient commutes with arbitrary base change.

\textbf{Proof.} Work étale locally on \(S=\operatorname{Spec}R\), with \(H=D_R(M)\). Its character group \(M\) is finitely generated. Right translation gives a grading \(A=\Gamma(G,\mathcal O_G)=\bigoplus_{m\in M}A_m\). Put \(B=A_0\) and \(Q=\operatorname{Spec}B\). Invariants commute with any base change because they are a direct grading summand. The same summand proves that \(B\) is flat over \(R\).

We check finite presentation. The data descend to a finitely generated \(\mathbf Z\)-algebra, including the smooth group, the closed embedding of \(D(M)\), and centrality; all are finite-presentation equations. We may therefore assume \(R\) Noetherian. Choose finitely many homogeneous generators of \(A\) as an \(R\)-algebra and a graded surjection \(P=R[z_1,\ldots,z_n]\to A\). Its invariant part surjects onto \(B\). The monomials of weight zero in \(P\) are finitely generated as a semigroup: their indecomposable exponent vectors form an antichain in \(\mathbf N^n\), since if \(u\le v\) are two such vectors, \(v-u\) also has weight zero. Every antichain is finite. The latter assertion follows inductively on \(n\) by extracting a nondecreasing subsequence of the first coordinates from any infinite sequence, then applying the induction to the remaining coordinates. Every zero-weight vector is a sum of indecomposables by induction on its coordinate sum. Thus \(P_0\), and then \(B\), are finitely generated. They are finitely presented over Noetherian \(R\). Descent and base change now give finite presentation in general.

Next we prove that \(q:G\to Q\) is an \(H\)-torsor. First suppose \(R=k\) is algebraically closed. The closed embedding \(H\hookrightarrow G\) gives a graded surjection \(A\to k[M]\). For each of finitely many generators \(m_i\) of \(M\), choose elements \(a_i\in A_{m_i}\) and \(b_i\in A_{-m_i}\) restricting to the corresponding characters of \(H\). Then \(f=\prod_i a_ib_i\in B\) has value one at the identity. In \(A_f\), every \(a_i\) and \(b_i\) is a unit. Each homogeneous component is therefore a free rank-one \(B_f\)-module generated by a homogeneous unit; multiplication of components is an isomorphism. Explicitly the map

\[
A_f\otimes_{B_f}A_f\longrightarrow A_f\otimes_k k[M],
\qquad a\otimes c_m\longmapsto ac_m\otimes e_m
\]

is an isomorphism: for a homogeneous unit \(u_m\) its inverse sends \(a\otimes e_m\) to \(a u_m^{-1}\otimes u_m\). Also \(A_f=\bigoplus_m B_f u_m\) is faithfully flat over \(B_f\). This proves the torsor assertion on \(D(f)\).

Left translation by \(g\in G(k)\) preserves the grading and gives the same proof on \(D(f_g)\), where \(f_g(x)=f(g^{-1}x)\). These opens cover \(Q\). Indeed \(B\hookrightarrow A\) has a \(B\)-linear retraction, so \(A\otimes_B\kappa(y)\) is nonzero for every prime \(y\) of \(B\) and \(q\) is surjective. A closed point of \(Q\) has a \(k\)-point \(g\) above it, and \(f_g(g)=1\). Since \(Q\) is of finite type over \(k\), an uncovered nonempty closed subset would contain a closed point. Hence there is no uncovered point.

Return to general \(R\). Both \(G\) and \(Q\) are flat of finite presentation over \(R\), and the geometric-fibre morphisms are faithfully flat torsors by the preceding calculation. The fibre criterion for flatness gives that \(q\) is flat and of finite presentation; it is surjective by the invariant retraction. The natural morphism \(G\times H\to G\times_QG\) is an isomorphism on geometric fibres. Both sides are flat of finite presentation over \(R\). The fibre criterion applied to this morphism makes it étale, and its fibrewise bijectivity makes it an isomorphism. Thus \(q\) is an fppf torsor. Centrality permits multiplication and inversion to descend, making \(Q\) a group; this also identifies \(Q\) with the fppf quotient.

Its geometric fibres are smooth connected reductive groups. Here is a proof of smoothness and reductivity that does not assume a quotient theorem. Over an algebraically closed field, \(A\) is a normal domain, and its invariant subring \(B\) is normal: an element of \(\operatorname{Frac}B\) integral over \(B\) is integral over \(A\), hence in \(A\), and is invariant, hence in \(B\). A finite-type normal group over a perfect field has a nonempty smooth open subset; translation carries a smooth point to every closed point, so it is smooth. It is connected because \(G\to Q\) is surjective. If its smooth connected unipotent radical were nontrivial, its inverse image would be a solvable normal subgroup of \(G\). The reduced identity component of that inverse image is smooth, connected, solvable and normal, so is a central torus by reductivity. It surjects onto the radical: the finitely many components of the inverse image cannot map onto a connected group with the identity component having smaller dimension. A torus has trivial image in a unipotent group, a contradiction. Flatness, finite presentation and these smooth geometric fibres prove smoothness of \(Q/S\).

\(\square\)

\section{2. Why rank one has only two moving directions}\label{AG-RG-03-why-rank-one-has-only-two-moving-directions}

We first prove the geometric rank-one statement, rather than assuming a root classification. Work over an algebraically closed field \(k\).

\textbf{Lemma 2.1.} If a torus acts linearly on an irreducible projective variety \(X\), it has at least \(\dim X+1\) fixed points.

\textbf{Proof.} Put \(d=\dim X\) and let \(T\) be the acting torus. Since \(k\) is algebraically closed, an étale splitting cover of \(T\) has a \(k\)-point: take a nonempty finite-type affine chart and apply the earlier Nullstellensatz. Pulling the splitting back along this point gives \(T\simeq\mathbf G_m^r\). The weight decomposition proved in the preceding lesson, \emph{Diagonal actions and their descent}, applies to the finite-dimensional representation defining the linear action. Choose homogeneous coordinates in which the action on points has the form \[
 t\cdot[c_0:\cdots:c_N]
       =[\chi_0(t)c_0:\cdots:\chi_N(t)c_N],
 \qquad \chi_j\in X^*(T)=\mathbf Z^r.
\] Delete any coordinates vanishing identically on \(X\), thereby restricting to an invariant projective linear subspace.

Here is an explicit cocharacter separating all distinct characters in this finite list. If \(r>0\), write \(\chi=(a_1,\ldots,a_r)\). For distinct characters \(\chi,\psi\), the integer polynomial \[
 \sum_{i=1}^r(a_i-b_i)M^{i-1}
\] is nonzero. The polynomial root bound proved in S02 Lemma 4.A, applied over \(\mathbf Q\), makes its set of integer roots finite. Choose a positive integer \(M\) outside the union of these finitely many sets and put \[
 \lambda(u)=(u,u^M,\ldots,u^{M^{r-1}}):
                  \mathbf G_m\longrightarrow T.
\] The integers \(w_j=\langle\chi_j,\lambda\rangle\) are equal only for equal characters. For \(r=0\), the unique cocharacter has this property since there is only the trivial character. Reorder the coordinates so that \(w_0\) is minimal. By the deletion step, \(c_0\) does not vanish identically on \(X\). The coordinate hyperplane \(H=V(c_0)\) is \(T\)-invariant.

We induct on \(d\). For \(d=0\), the reduced irreducible finite-type \(X\) has one point. An affine neighborhood is a zero-dimensional domain, hence a field; the Nullstellensatz makes it \(k\). Thus \(X=\operatorname{Spec}k\), and its point is fixed.

Suppose \(d>0\). The section \(X\cap H\) is nonempty. Otherwise \(X\) would be closed in the affine complement of \(H\) and hence affine. S01 Theorem 8.4 and its closed-immersion cohomology calculation prove that \(H^0(X,\mathcal O_X)\) is finite-dimensional over \(k\). It is a domain because \(X\) is integral. Multiplication by a nonzero element is an injective endomorphism of a finite-dimensional vector space and therefore surjective, so this algebra is a field. Algebraic closedness makes it \(k\). Affineness would then give \(X=\operatorname{Spec}k\), contradicting \(d>0\).

Every irreducible component of \(X\cap H\), with its reduced structure, has dimension \(d-1\). To check the dimension rather than assume it, take a standard projective affine chart meeting the component and intersect it with \(X\). Its algebra \(A\) is a finite-type domain of dimension \(d\), by the function-field dimension formula in AG-CA \emph{Krull dimension and Noether normalization}, Theorem 4.2. In a local frame the hyperplane section is defined by a nonzero element \(f\in A\); its component corresponds to a prime \(\mathfrak p\) minimal over \((f)\). This prime is nonzero. The principal ideal theorem in AG-CA \emph{Dimension theory of Noetherian local rings}, Theorem 3.1, gives \(\operatorname{ht}\mathfrak p\le1\), while \((0)\subsetneq\mathfrak p\) gives the reverse inequality. The height formula, Theorem 4.3 of the former lesson, yields \[
 \dim(A/\mathfrak p)=\dim A-\operatorname{ht}\mathfrak p=d-1.
\] Passing between nonempty affine opens of this integral component preserves its function field and hence its dimension by that same Theorem 4.2.

Choose one such reduced component \(Y\). It is \(T\)-stable. Indeed, \(T\times Y\) is integral: on affine charts of \(Y\) its algebra is the domain \(A[t_1^{\pm1},\ldots,t_r^{\pm1}]\), and these charts have nonempty intersections. The closure of its image under the action is therefore irreducible, lies in \(X\cap H\), and contains \(Y\) through the identity of \(T\). Maximality of the irreducible component forces this closure to be \(Y\). The action factors through its reduced structure, since pullbacks of its defining ideal vanish at every point of the reduced source and therefore vanish. The restriction is a linear action on the projective variety \(Y\). Induction supplies at least \(d\) distinct \(T\)-fixed points in \(Y\subset H\).

There is one more fixed point outside \(H\). Choose \(a\in(X\setminus H)(k)\), using the Nullstellensatz on this nonempty affine chart, and normalize its coordinates so that \(c_0(a)=1\). Its \(\lambda\)-orbit on the same chart has coordinates \[
 c_j(\lambda(u)a)/c_0(\lambda(u)a)
       =u^{w_j-w_0}c_j(a),\qquad j>0.
\] All exponents are nonnegative integers, so these polynomials define a morphism \(\mathbf A^1_k\to\mathbf P^N_k\setminus H\). It lands in \(X\): the equations of the closed affine subscheme \(X\setminus H\) vanish after substitution in \(k[u,u^{-1}]\), hence in its subring \(k[u]\). At \(u=0\), every coordinate of weight greater than \(w_0\) is zero; the remaining coordinates have character \(\chi_0\) by the separating property. Every element of \(T\), over every \(k\)-algebra, scales these remaining coordinates by the same unit \(\chi_0(t)\). Thus the resulting projective point is \(T\)-fixed. Its zeroth coordinate is still one, so it lies outside \(H\) and differs from all \(d\) points obtained inside \(H\). This proves the bound \(d+1\). \(\square\)

\subsection{The curve duality used in rank one}\label{AG-RG-03-the-curve-duality-used-in-rank-one}

\textbf{Lemma (curve duality).} If \(C\) is a smooth projective pure one-dimensional curve over a field \(k\) and \(L\) is invertible, then naturally

\[
H^1(C,L)^*\simeq
H^0(C,\Omega_{C/k}\otimes L^{-1}).
\]

\textbf{Proof.} Embed \(i:C\hookrightarrow P=\mathbf P^N_k\), and put \(c=N-1\). The regular-immersion calculation in \href{https://kokunoyumeto.github.io/open-math-courses/courses/AG-RG/AG-RG-S02.html}{Projective cohomology and smooth affine models}, Lemma 6.1, proves over the original field, including imperfect fields, that \(I/I^2\) is locally free of rank \(c\), that the ideal is locally generated by a regular sequence, and that \[
0\longrightarrow I/I^2\longrightarrow i^*\Omega_{P/k}
\longrightarrow\Omega_{C/k}\longrightarrow0
\] is exact. Its proof first calculates at geometric closed points, then descends the conormal module and the exact Koszul complex by faithful field extension.

Dualizing that Koszul resolution, with a local frame of \(L\), gives \[
\mathcal E xt^q_P(i_*L,\omega_P)=0\quad(q\ne c),
\qquad
\mathcal E xt^c_P(i_*L,\omega_P)
=i_*(L^{-1}\otimes i^*\omega_P\otimes\det(I/I^2)^*).
\] Here \(\omega_P=\det\Omega_{P/k}=\mathcal O_P(-N-1)\). Taking determinants of the conormal sequence identifies the sheaf in degree \(c\) with \(i_*(L^{-1}\otimes\Omega_{C/k})\). The finite-cover Hom double complex written in S02 Theorem 6.2 has only this nonzero sheaf-Ext row, and therefore gives \[
\operatorname{Ext}^c_P(i_*L,\omega_P)
=H^0(C,L^{-1}\otimes\Omega_{C/k}).
\] S02 Ambient Theorem 6.B proves projective-space duality by the monomial trace, a two-stage twist presentation and dimension shifting; it identifies the same Ext group with \(H^{N-c}(P,i_*L)^*=H^1(C,L)^*\). Its trace sends the class of \(T_0^{-1}\cdots T_N^{-1}\) to one. All the Koszul, determinant, finite-cover and trace comparisons are natural. This proves the asserted isomorphism. \(\square\)

\textbf{Lemma 2.2 (the curve argument).} Over an algebraically closed field, a smooth projective curve dominated by \(\mathbf P^1\) is isomorphic to \(\mathbf P^1\).

\textbf{Proof.} Write \(C\) for the dominated curve. The dense image of the irreducible \(\mathbf P^1_k\) is irreducible. Smoothness gives reducedness, so \(C\) is integral and connected. If domination is initially given by a rational map, the homogeneous-coordinate DVR extension in the earlier S03 Theorem 2.2 makes it a morphism. Thus the hypotheses of \href{https://kokunoyumeto.github.io/open-math-courses/courses/AG-RG/AG-RG-S02.html}{S02 Proposition 6.5} hold: its complete separable differential calculation and inseparable-subfield reduction prove \(g=\dim_kH^1(C,\mathcal O_C)=0\) in every characteristic. Its proof is earlier in the programme; no unproved external form of Lüroth's theorem is used.

We supply explicitly the concluding rational-function calculation from that proposition. Choose a closed point \(q\in C(k)\), which exists by the earlier Nullstellensatz. S02 Lemma 6.3 proves \(H^0(C,\mathcal O_C)=k\), divisor Euler characteristics, the canonical degree \(\deg\Omega_{C/k}=2g-2=-2\), and vanishing for negative-degree line bundles. Its earlier Theorem 6.2 gives duality. For every integer \(m\ge0\), these results yield \[
 H^1(C,\mathcal O_C(mq))^*
   =H^0(C,\Omega_{C/k}(-mq))=0,
 \qquad
 h^0(C,\mathcal O_C(mq))=m+1.
\] The first vanishing uses degree \(-2-m<0\); the second is the proved Euler characteristic formula with \(g=0\).

Take \(x\) in \(H^0(C,\mathcal O_C(q))\) independent of the constants. In the divisor description of this sheaf, it is a rational function regular away from \(q\) with pole order at most one at \(q\). It has a pole: otherwise it would be a global regular function and belong to \(k\). Its pole order is therefore exactly one. The functions \(1,x,\ldots,x^m\) belong to \(H^0(C,\mathcal O_C(mq))\) and have distinct pole orders. A nonzero linear combination has the pole order of its highest nonzero power, since all coefficients lie in \(k\) and that term cannot be canceled by terms of smaller order. They are independent, and their number \(m+1\) is the dimension just computed. They form a basis.

Let \(h\in k(C)\) be any nonzero rational function. The earlier divisor proof shows that it has finitely many poles. List its poles \(a_1,\ldots,a_r\) outside \(q\). The function \(x\) is regular at each of them, and \(x(a_i)\in k\). Each nonzero function \(x-x(a_i)\) has positive valuation at \(a_i\) and no pole outside \(q\). Choose integers \(N_i\ge0\) large enough that \[
 P(x)h,\qquad P(T)=\prod_{i=1}^r(T-x(a_i))^{N_i},
\] has no pole outside \(q\). Multiplying cannot introduce a pole at another such point. Its remaining pole order at \(q\) is finite, so the product is a section of \(\mathcal O_C(mq)\) for some \(m\ge0\). By the basis calculation, it equals a polynomial \(Q(x)\). The polynomial \(P\) is nonzero and \(x\) is transcendental over the algebraically closed \(k\), so \(P(x)\ne0\). Consequently \(h=Q(x)/P(x)\in k(x)\). The empty pole list gives \(P=1\), and the zero function is already in \(k(x)\). Hence \(k(C)=k(x)\).

This equality gives inverse rational maps between \(C\) and \(\mathbf P^1_k\): on a nonempty affine chart, send its finitely many coordinate generators to their rational functions in the other function field and clear their denominators; their equations remain satisfied. Both targets are projective and both sources are smooth curves. The earlier S03 Theorem 2.2 extends these maps at every DVR by homogeneous coordinates. Their composites are the identity on the common function field and hence on a dense open. The closed-diagonal equalizer argument for the reduced integral sources makes the composites the identity everywhere. The resulting morphisms are inverse isomorphisms. This argument needs no degree formula for an unchecked inseparable map. \(\square\)

\textbf{Proposition 2.3.} A semisimple group of rank one has exactly two Borel subgroups containing a maximal torus, its Borel quotient is \(\mathbf P^1\), and the action on that quotient has image \(\operatorname{PGL}_2\) and kernel its centre.

\textbf{Proof.} Fix a maximal torus \(T\subset B\) and put \(X=G/B\). RG01 Lemma 2.1 supplies its linear equivariant projective embedding; the earlier quotient-smoothness proof in AG-GS, \emph{Quotients and torsors}, Theorem 11.7 makes \(X\) smooth. It is integral, being reduced and the surjective image of the irreducible smooth connected \(G\). RG02's complete torus-fixed-coset proof makes \(X^T\) a finite disjoint union of reduced \(k\)-points. Its following transitivity argument identifies this set with \(N_G(T)/C_G(T)\). RG02 Theorem 2.1 gives \(C_G(T)=T\), so this quotient acts faithfully on the rank-one character lattice \(\mathbf Z\), whose automorphisms are precisely multiplication by \(1\) or \(-1\). Hence \(X^T\) has at most two points. It has at least one, the coset \(B\). If \(\dim X=0\), integrality and the Nullstellensatz make \(X=\operatorname{Spec}k\), so the torsor stabilizer forces \(G=B\), contradicting nontrivial semisimplicity. Lemma 2.1 now gives \(|X^T|\ge\dim X+1\ge2\). Therefore there are exactly two fixed points and \(\dim X=1\). The Borel-normalizer theorem already proved in RG02 identifies these with exactly two Borels containing \(T\).

Let \(K\) be the schematic kernel of the action on \(X\). It is the intersection of the conjugates of \(B\). This retains the full kernel on every test algebra: in the equivariant projective embedding the fixed-point equations are the coordinate two-by-two minors. If they vanish on every constant \(k\)-point of the reduced \(X\), each of their finitely many linearly independent scalar coefficients vanishes on every such point, hence vanishes identically by the Nullstellensatz. Thus these equations vanish after arbitrary scalar extension as well. Its reduced identity component is a smooth connected normal solvable subgroup, hence is trivial by semisimplicity. Normality here holds schematically: conjugation by the smooth connected \(G\) preserves that component, and its reduced source factors through its reduced structure, as in S08 Lemma G.4.1. Thus \(K\) is zero-dimensional. It is finite: its affine algebra is Noetherian, has finitely many minimal primes, and in dimension zero they are maximal, with residue fields \(k\). The nilradical is generated by finitely many nilpotent elements and hence is nilpotent, by the bound on monomial exponents. The reduced algebra is a finite product of copies of \(k\); the successive nilradical quotients are finite modules over it. This gives a finite-dimensional algebra for \(K\). A positive-dimensional Borel consequently has a nonconstant orbit map to \(X\). Otherwise all its point orbits would be points, and density of rational points in the reduced \(B\times X\) would make the action trivial schematically, contradicting \(B\subset K\).

We supply the rational coordinates required for restricting such an orbit map to a line. For a smooth connected solvable group \(B=U\rtimes T\), RG01 Lemma 2.B gives \[
 U=U_0\supset U_1\supset\cdots\supset U_a=1,
 \qquad U_i/U_{i+1}\simeq\mathbf G_a.
\] The quotients \(U/U_i\) are affine by the earlier normal affine-quotient proof, AG-GS, \emph{Quotients and torsors}, Theorem 11.1d. Its Theorem 11.16 makes \[
 U/U_{i+1}\longrightarrow U/U_i
\] a \(\mathbf G_a\)-torsor. That theorem uses left cosets; here all \(U_i\) are normal in \(U\), so left and right cosets give the same quotient group and the same assertion. Every such torsor over an affine algebra has a section by the explicit vector-line module-descent calculation in RG01's solvable-structure proof: the translation cocycle descends \(0\to A\to E\to A\to0\), and a lift of \(1\) gives the section. This calculation works over any affine \(A\). Translation by a section identifies the torsor with its base times \(\mathbf A^1\). Starting with \(U/U_0=1\) and repeating gives \(U\simeq\mathbf A^a\) as varieties. The torus is split over the algebraically closed field, by the splitting argument in Lemma 2.1. Thus \[
 B\simeq\mathbf A^a\times\mathbf G_m^r,
 \qquad k(B)=k(z_1,\ldots,z_d),\quad d=a+r.
\] These are scheme coordinates; no assertion about abstract group coordinates is needed.

Choose a nonconstant orbit map \(b\mapsto bx\) and a nonconstant \(h\in k(X)\). The map is dominant because its image has positive dimension in the integral curve, so its pullback \(f\in k(B)\) is nonconstant. Choose the least \(j\ge1\) for which \(f\in k(z_1,\ldots,z_j)\). Then \(f\notin k(z_1,\ldots,z_{j-1})\). Write, after clearing finitely many coefficient denominators, \[
 f=P(z_j)/Q(z_j),\quad
 P=\sum_i A_i z_j^i,\quad Q=\sum_i B_i z_j^i,
 \quad A_i,B_i\in k[z_1,\ldots,z_{j-1}].
\] Some coefficient minor \(A_iB_l-A_lB_i\) is nonzero. If they all vanished, choosing \(B_l\ne0\) would give \(P=(A_l/B_l)Q\), contradicting the choice of \(j\).

The inverse image of the domain where \(h\) is regular contains a nonempty principal open of our coordinate chart. Choose a nonzero polynomial \(F\) defining such an open, multiplying by torus coordinates to clear Laurent denominators if necessary. Viewed as a polynomial in \(z_j\), it has a nonzero coefficient in the other variables. Specialize every other coordinate so that this coefficient, the selected minor, and all fixed torus coordinates stay nonzero. Such a specialization exists: their product is a nonzero polynomial over the infinite field \(k\), and repeated use of the univariate root bound in S02 Lemma 4.A supplies a point where it does not vanish. For no remaining variables the product is already a nonzero scalar. On the resulting coordinate line, \(P/Q\) stays nonconstant and \(F\) does not vanish identically. If \(z_j\) is a torus coordinate, restrict also to \(z_j\ne0\). We have obtained a nonconstant rational map \(\mathbf A^1\dashrightarrow X\), since its composition with \(h\) is this nonconstant rational function. The explicit DVR extension calculation in S03 Theorem 2.2 extends it to \(\mathbf P^1\to X\). Lemma 2.2, which includes inseparable domination, identifies \(X\) with \(\mathbf P^1\).

The natural map \(\operatorname{PGL}_2\to\operatorname{Aut}(\mathbf P^1)\) is an isomorphism on all test algebras. An automorphism sends \(0,1,\infty\) to three sections that remain pairwise disjoint on every fibre. Locally trivialize their quotient lines. A change of basis sends \(0,\infty\) to the first and third; the two nonzero coordinates of the middle section are units, so diagonal scaling also sends \(1\) to it. For uniqueness, an automorphism fixing \(0\) and \(\infty\) preserves both Cartier divisors. The ratio of its pulled-back coordinate to that coordinate is therefore an invertible regular function on all of \(\mathbf P^1\). S01 Theorem 7.2 gives \(H^0(\mathbf P^1_R,\mathcal O)=R\), so that ratio is a base unit. Fixing \(1\) makes it one. The local projective matrices consequently agree on overlaps and glue in the represented projective linear group. This proves the assertion even over rings with nilpotents.

The schematic image \(H\) of \(G\to\operatorname{PGL}_2\) is smooth and connected, and the map onto it is faithfully flat, by AG-GS, \emph{Group schemes over a field}, Proposition 5.8 and the perfect-field smoothness proof. Its kernel \(K\) is finite. Its torus image has dimension one, and maximal-torus conjugacy in RG01 conjugates it to the diagonal torus of \(\operatorname{PGL}_2\). If \(H\) were proper, it would have dimension at most two. In every characteristic, \(\operatorname{Lie}\operatorname{PGL}_2=\mathfrak{gl}_2/kI\) has the diagonal line of weight zero and two distinct root lines of weights \(1,-1\) for the cocharacter represented by \(\operatorname{diag}(t,1)\). The \(T\)-stable Lie algebra of \(H\) contains the diagonal line and at most one root line. Choose that cocharacter or its inverse so that all its weights on \(\operatorname{Lie}H\) are nonnegative. The general limit-subgroup proof already given in RG01 Lemma 3.3 makes \(P_H(\lambda)\) smooth with this entire Lie algebra, hence of dimension \(\dim H\). A closed subgroup of that dimension in the connected smooth \(H\) is all of \(H\). But \[
 P_H(\lambda)=H\cap P_{\operatorname{PGL}_2}(\lambda),
\] and the ambient limit subgroup is triangular. Thus \(H\) would be solvable. The finite group \(K(k)\) is central in \(G(k)\): each conjugation orbit is the image of the connected reduced \(G\) in a finite set and is the identity translate. Every \(H(k)\)-point lifts through the faithfully flat finite-type image map by the Nullstellensatz. Hence \(G(k)\) is a central extension of a solvable group and is solvable. Its finite derived-word identities hold on the dense rational points of each reduced finite Cartesian power of \(G\), so they hold schematically, making \(G\) solvable. This contradicts semisimplicity. Therefore \(H=\operatorname{PGL}_2\).

In particular \(\dim G=3\). Since \(\dim X=1\), each of the two Borels \(B_+,B_-\) has dimension two and unipotent radical of dimension one. The one-dimensional unipotent classification proved in S03 Section 4 identifies each radical with \(\mathbf G_a\). Its image has dimension one because the kernel is finite. Projection of the corresponding point stabilizer onto its diagonal torus gives a character of the radical, trivial by RG01 Lemma 2.A. Thus its reduced closed image lies in the ambient unipotent line; dimension one makes it that entire line. Its torus character \(\alpha_+\) or \(\alpha_-\) is nonzero, since the induced diagonal-torus map is an isogeny and the ambient root character is nonzero. The normalizer element interchanging the two fixed points acts by inversion on \(T\), so \(\alpha_-=-\alpha_+\).

The Lie algebra of \(K\) lies in both \(\mathfrak b_+\) and \(\mathfrak b_-\). Their nonzero-weight lines have these distinct opposite characters, while the shared torus has weight zero. Thus \(\operatorname{Lie}K\) has no vector in either unipotent root line. Each finite intersection \(K\cap R_u(B_\pm)\) is torus-stable. Since \(\alpha_\pm\ne0\), its polynomial ideal is homogeneous for the distinct weights \(i\alpha_\pm\), hence is \((x^n)\) for some \(n\ge1\). The Hopf identity requires every intermediate binomial coefficient in \((X+Y)^n\) to vanish. In characteristic zero the coefficient \(n\) excludes \(n>1\). In characteristic \(p\), write \(n=p^e m\) with \(p\nmid m\). If \(m>1\), the coefficient \(m\) of \(X^{p^e}Y^{n-p^e}\) in \((X^{p^e}+Y^{p^e})^m\) is nonzero, a contradiction. Thus \(n=p^e\). Any \(n>1\) gives nonzero tangent space at the identity, already excluded. Both intersections are therefore trivial. The image maps of the two radicals are faithfully flat with trivial kernel, hence isomorphisms onto the ambient unipotent subgroups, by their full torsor relations.

For a point of \(B_+\cap B_-\) over any test algebra, write it \(tu_+\) by the solvable decomposition. Its image lies in the diagonal intersection of the two ambient point stabilizers, so the image of \(u_+\) is trivial. The just-proved isomorphism gives \(u_+=1\). Therefore \(B_+\cap B_-=T\) schematically and \(K\subset T\). A finite normal multiplicative-type subgroup of a smooth connected group is central: over \(k\) it has a finite character group, and the earlier coefficient proof for diagonalizable groups represents its automorphisms by that finite constant étale character-automorphism group. The conjugation map from connected \(G\) to this discrete group has value the identity automorphism at the identity of \(G\), hence is constant with that value as a scheme morphism. Thus \(K\subset Z(G)\). Conversely \(Z(G)\) maps to the centre of \(\operatorname{PGL}_2\), which is trivial schematically. Indeed a central automorphism of \(\mathbf P^1\) commutes with the diagonal torus and preserves its fixed sections \(0,\infty\); interchanging them would invert the universal torus parameter and cannot commute with it. It is therefore a scaling by a base unit, and commuting with every translation forces that unit to be one. Thus \(Z(G)\subset K\), proving \(K=Z(G)\) and every assertion. \(\square\)

For a reductive group of semisimple rank one, quotienting by its radical gives this situation. Its centre is the kernel of the resulting action on the same projective line. The two Borel unipotent radicals are still additive lines with opposite characters.

Now let \(\alpha\) be any nonzero weight of a maximal torus of an arbitrary reductive \(k\)-group. Define the codimension-one torus

\[
T_\alpha=(\ker\alpha)^0_{\mathrm{red}}.
\]

The centralizer \(G_\alpha=C_G(T_\alpha)\) is reductive by the preceding lesson. Its weights are precisely the weights of \(G\) that are rational multiples of \(\alpha\), since these and only these vanish on \(T_\alpha\). Its maximal central torus has codimension one in \(T\): the nonzero characters on \(T/T_\alpha\) span a rank-one lattice. Thus its semisimple rank is one. Proposition 2.3 proves

\[
\Phi\cap\mathbf Q\alpha=\{\alpha,-\alpha\},\qquad
\dim\mathfrak g_\alpha=1.
\]

This proves both the reducedness along each root line and the one-dimensional root-space assertion without appealing to a general root-system classification.

\section{3. Subgroups obtained by taking a limit}\label{AG-RG-03-subgroups-obtained-by-taking-a-limit}

The definitions and full grading, smoothness, connectedness and open-cell proof for these subgroups are in \href{https://kokunoyumeto.github.io/open-math-courses/courses/AG-RG/AG-RG-01.html\#subgroups-obtained-by-taking-a-limit}{Tori, maximal tori and their conjugacy, Lemma 3.3}, an earlier lesson. In this lesson Lemma 3.1 denotes that already proved general lemma; it applies to every smooth affine group of finite presentation, without reductivity or a root classification.

For \(\operatorname{SL}_2\) and \(\lambda(t)=\operatorname{diag}(t,t^{-1})\), conjugation sends \(\begin{pmatrix}a&b\\c&d\end{pmatrix}\) to \(\begin{pmatrix}a&t^2b\\t^{-2}c&d\end{pmatrix}\). Thus \(P(\lambda)\) is upper triangular, \(U(\lambda)\) is upper unipotent, and \(L(\lambda)\) is diagonal. This calculation works in characteristic two because the group character \(t^2\) remains nontrivial.

\section{4. Root groups over the base}\label{AG-RG-03-root-groups-over-the-base}

Let \(T=D_S(M)\) be a split maximal torus. Locally the nonzero weights are a fixed finite set \(\Phi\subset M\), and each \(\mathfrak g_\alpha\) is a line bundle by Proposition 2.3 on geometric fibres. Define \(T_\alpha\) by the torsion-free quotient of \(M/\mathbf Z\alpha\). Then

\[
G_\alpha=C_G(T_\alpha),\qquad
\operatorname{Lie}(G_\alpha)=\mathfrak t\oplus\mathfrak g_\alpha\oplus\mathfrak g_{-\alpha}.
\]

Choose a cocharacter \(\lambda\) of \(T\) with \(n=\langle\alpha,\lambda\rangle>0\), and set \(U_\alpha=U_{G_\alpha}(\lambda)\). It is a smooth connected one-dimensional unipotent group, with Lie algebra \(\mathfrak g_\alpha\).

For a line bundle \(L\), write \(\mathbf V(L)=\operatorname{Spec}_S\operatorname{Sym}(L^*)\) for its additive group.

\textbf{Theorem 4.1.} There is a unique \(T\)-equivariant homomorphism

\[
\exp_\alpha:\mathbf V(\mathfrak g_\alpha)\longrightarrow G
\]

whose differential is the inclusion of the root line. It is a closed immersion, with image \(U_\alpha\), and commutes with every base change. Moreover

\[
U_{-\alpha}\times T\times U_\alpha\longrightarrow G_\alpha
\]

is an open immersion.

\textbf{Proof.} On \(U_\alpha\), \(\mathbf G_m\) acts through \(t^n\) on each geometric fibre. Its subgroup \(\mu_n\) acts trivially: the fixed subgroup is smooth by exactness of multiplicative-type invariants, and equals \(U_\alpha\) on every geometric fibre, so equals it over \(S\).

Étale locally choose a section \(\sigma\) of \(U_\alpha\) disjoint from its identity. The orbit map \(t\mapsto t.\sigma\) extends to \(q:\mathbf A^1\to U_\alpha\) with \(q(0)=1\). Since \(\mu_n\) acts trivially on the target, the map factors through the invariant ring \(R[x]^{\mu_n}=R[x^n]\) on an affine chart. Its factor \(\bar q:\mathbf A^1\to U_\alpha\) is an isomorphism on geometric fibres: there the original map is a nonzero scalar times \(t^n\). The smooth fibrewise isomorphism criterion makes \(\bar q\) an isomorphism over the base.

The quotient acting torus \(\mathbf G_m/\mu_n\) gives ordinary scalar multiplication on the source. The transported group law is a polynomial \(m(x,y)\) with

\[
m(tx,ty)=tm(x,y),\quad m(x,0)=x,\quad m(0,y)=y.
\]

Homogeneity forces \(m=ax+by\), and the identities force \(a=b=1\). Thus \(U_\alpha\) is an additive line. Normalize the isomorphism by its differential to obtain \(\exp_\alpha\). A scalar-equivariant polynomial map of an additive line is linear; if its differential is the identity it is the identity. This proves uniqueness and allows descent of the local construction.

An equivariant homomorphism with the required differential must factor through \(G_\alpha\), since \(T_\alpha\) acts trivially on its source, and through \(U_{G_\alpha}(\lambda)\), since its orbit limits to the identity. The same uniqueness therefore holds as a map into \(G\). The open immersion follows from Lemma 3.1, with \(L_{G_\alpha}(\lambda)=T\). \(\square\)

The notation \(\exp_\alpha\) denotes this algebraic map, not the analytic exponential series. No division by factorials occurs. A root-group parameterization \(x_\alpha:\mathbf G_a\simeq U_\alpha\) is exactly a trivializing section of the root line.

\textbf{Example 4.2.} In \(\operatorname{GL}_n\), conjugation by \(\operatorname{diag}(t_i)\) sends \(E_{ij}\) to \((t_i/t_j)E_{ij}\). Thus

\[
\alpha_{ij}=e_i-e_j,\quad
\mathfrak g_{\alpha_{ij}}=\mathcal O_SE_{ij},\quad
x_{ij}(u)=I+uE_{ij}.
\]

\textbf{Example 4.3.} For a line bundle \(L\) on \(S\), consider \(G=\operatorname{SL}(\mathcal O_S\oplus L)\). It has the split torus acting by \((t,t^{-1})\) on the two summands. Its two root lines are \(L^{-1}\) and \(L\). If \(L\) is nontrivial, its root groups are \(\mathbf V(L^{-1})\) and \(\mathbf V(L)\) and need not be globally isomorphic to \(\mathbf G_a\). Having a split maximal torus is thus weaker than having globally parameterized root groups. The distinction disappears over a field.

\section{5. A projective quotient and its rank-one form}\label{AG-RG-03-a-projective-quotient-and-its-rank-one-form}

We use the general algebraic-space tool proved before this lesson in \href{https://kokunoyumeto.github.io/open-math-courses/courses/AG-RG/AG-RG-S04.html}{Affine descent, Zariski Main and recognition of spaces}, Theorems E8.1 and E9.1. The proof there includes the finite flat affine quotient, the off-diagonal induction, separated locally quasi-finite scheme effectivity, and relative integral closure with its étale base change. It applies over every base and on every test scheme.

\textbf{Quasi-finite normalization lemma.} Let \(S\) be a scheme and let \(f:X\to Y\) be a quasi-finite separated morphism of algebraic spaces over \(S\). Let \(Y'\) be the relative integral closure of \(Y\) in \(X\). Then the factorization \(X\xrightarrow{f'}Y'\xrightarrow{\nu}Y\) has \(f'\) a quasi-compact open immersion and \(\nu\) integral. In particular \(f\) is quasi-affine.

\textbf{Proof.} Theorem E9.1 of the preceding supporting lesson proves the full factorization, after Theorem E8.1 establishes recognition and Corollary E8.2 establishes representability on every scheme test. Its integral algebras are descended using the proved affine faithfully flat equalizer, and Corollary D4.3 supplies their étale base-change compatibility. Thus for every affine scheme \(Z\to Y\), the pullback \(Z\times_YX\) is a quasi-compact open subscheme of the affine integral scheme \(Z\times_YY'\). This is quasi-affineness. All these input proofs precede the present application in the programme. \(\square\)

The relative integral closure retains nilpotents; a reduced normalization is not substituted for it. The full argument is in the preceding supporting lesson, rather than being delegated to a source citation.

The following construction will also be used for flag varieties. It makes explicit why a fibrewise projective quotient is a scheme, even before we have classified the group over the base.

\textbf{A flat-source fibre criterion.} A finitely presented closed immersion \(Y\hookrightarrow Z\) over \(S\), with \(Y\) flat over \(S\), is an isomorphism if it is so on every geometric fibre. Locally write its ideal as \(I\subset C\). Flatness of \(C/I\) makes \(I\otimes k(s)\to C\otimes k(s)\) injective. Fibre equality makes this module zero. Since \(I\) is finitely generated as an ideal, Nakayama at every prime of \(C\) gives \(I=0\). No flatness assumption on \(Z\) is needed.

\subsection{Formal separation and normalizers}\label{AG-RG-03-formal-separation-and-normalizers}

We first justify the use of identity neighbourhoods over bases with nilpotents. This argument concerns smooth connected subgroups, rather than only tori.

\textbf{Formal separation lemma.} Let \(R\) be Noetherian and let \(H/R\) be a smooth affine group of finite presentation with geometrically connected fibres. Write \(C=\mathcal O(H)\) and let \(J\) be its augmentation ideal. Then

\[
\bigcap_{n\geq 1}J^n=0.
\]

\textbf{Proof.} A smooth geometrically connected algebraic group over a field is geometrically integral: over an algebraic closure its regular irreducible components are disjoint and open, and connectedness leaves one component. For every prime \(\mathfrak p\subset R\), the flat algebra \(C/\mathfrak p C\) therefore embeds in its integral generic fibre over \(\operatorname{Frac}(R/\mathfrak p)\) and is a domain.

Choose a finite prime filtration of the \(R\)-module \(R\). Such a filtration exists for every finite module over a Noetherian ring: choose an element with prime annihilator, insert its cyclic submodule, and repeat in the quotient; the ascending chain condition terminates the process. Flatness of \(C\) turns this into a filtration of \(C\) with factors \(C/\mathfrak p C\). If \(j\in J\), then \(1-j\) has augmentation one in every such factor, so is nonzero there. It is a nonzerodivisor on every factor and consequently on \(C\).

Put \(I=\bigcap_nJ^n\). The ring \(C\) is Noetherian. Artin--Rees applied to \(I\subset C\) gives, for some \(c\) and every \(n\geq c\),

\[
I\cap J^n=J^{n-c}(I\cap J^c).
\]

For completeness, this equality comes from the finite generation of the graded submodule \(\bigoplus_n(I\cap J^n)t^n\) of the Rees ring \(\bigoplus_nJ^nt^n\): the latter is a finitely generated \(C\)-algebra, hence Noetherian; take \(c\) at least the degrees of homogeneous generators. Since \(I\subset J^n\) for all \(n\), taking \(n=c+1\) gives \(I=JI\). On a finite set of generators of \(I\), this equality has matrix form \(v=Mv\) with entries of \(M\) in \(J\). The adjugate identity gives \(\det(1-M)I=0\), where \(\det(1-M)=1-j\) for some \(j\in J\). Its nonzerodivisor property proves \(I=0\). \(\square\)

\textbf{Normalizer lemma.} If \(H\) is a smooth closed subgroup with geometrically connected fibres of a smooth affine group \(K\) of finite presentation over any scheme \(S\), then its scheme-theoretic normalizer is a closed subgroup of finite presentation in \(K\). Its formation commutes with every base change.

\textbf{Proof.} First let \(S=\operatorname{Spec}R\) be Noetherian. The identity neighbourhood \(H_n\) has coordinate algebra \(C/J^{n+1}\), finite locally free over \(R\). Indeed, smooth coordinates along the identity section identify the augmentation graded algebra with \(\operatorname{Sym}(J/J^2)\); each of its finitely many graded pieces is finite locally free, and their extensions split locally. The same holds for \(K_n\), and \(H_n\subset K_n\) is a closed subscheme.

Write \(D=\mathcal O(K)\). Restrict the conjugates of a finite set of generators of the ideal of \(H\) to \(H_n\). Their vanishing gives finitely many coefficient equations in \(D\), because \(\mathcal O(H_n)\) is finite projective. Impose these equations for conjugation by \(g\) and by \(g^{-1}\). They define a closed subgroup \(N_n\subset K\), the stabilizer of \(H_n\). Let \(L_n\subset D\) be its ideal. The ideals \(L_n\) increase, so their union \(L\) equals some \(L_N\).

It remains to check that these finite equations detect the whole subgroup on every test ring. Over the universal quotient \(D/L\), both conjugation restrictions vanish on every identity neighbourhood of \(H\). This quotient is Noetherian, and \(H_{D/L}\) is still smooth with geometrically connected fibres. The formal separation lemma makes each restricted defining equation zero in \(\mathcal O(H_{D/L})\). Thus both conjugations preserve \(H_{D/L}\), and \(N_N\) normalizes \(H\). Pulling back this actual equality proves sufficiency over every \(R\)-algebra, including non-Noetherian ones. Necessity follows because a group automorphism preserving \(H\) preserves its identity and all its identity neighbourhoods. Therefore \(N_K(H)=N_N\) as functors, and the description proves arbitrary base change.

For general affine \(S\), \href{https://kokunoyumeto.github.io/open-math-courses/courses/AG-RG/AG-RG-S02.html}{Projective cohomology and smooth affine models}, Proposition 3.2 and Theorem 4.5, supply smooth affine group and subgroup models with geometrically connected fibres over a finitely generated \(\mathbf Z\)-algebra. Their proof descends the finite smoothness witnesses, proves constructibility of the number of geometric components, and removes the constructible bad locus at a finite coefficient stage. Apply that argument simultaneously to \(K,H\), retaining the closed immersion, group laws and their finitely many identities. The Noetherian construction above then represents the normalizer on every algebra over the model. Its pullback represents the original normalizer on every \(S\)-scheme, including non-Noetherian and nonreduced tests. Uniqueness of the represented subgroup glues these constructions. \(\square\)

\subsection{Cohomology of a projective completion}\label{AG-RG-03-cohomology-of-a-projective-completion}

The preceding supporting lesson \href{https://kokunoyumeto.github.io/open-math-courses/courses/AG-RG/AG-RG-S02.html}{Projective cohomology and smooth affine models}, Theorem 2.3, proves the precise formal-functions statement used here. If \(A\) is Noetherian, \(X/A\) projective, \(F\) coherent and \(I\subset A\), then \[
H^q(X,F)^\wedge_I
\simeq\varprojlim_nH^q(X,F/I^nF).
\] Here is the mechanism of that written proof. Over the Rees algebra \(B=\bigoplus I^nt^n\), the sheaf \(\bigoplus I^nF\) is finite on the projective base change \(X_B\); projective Serre finiteness therefore makes \(\bigoplus H^q(X,I^nF)t^n\) finite over \(B\). If \(M^q=H^q(X,F)\), \(J_n^q=\operatorname{im}(H^q(X,I^nF)\to M^q)\) and \(K_n^q\) is its kernel, finite homogeneous generators give constants \(a_q,b_q\) with \[
I^nM^q\subset J_n^q\subset I^{n-a_q}M^q,
\qquad K_m^q\longrightarrow K_n^q=0\quad(m\ge n+b_q).
\] The sheaf exact sequence consequently gives \[
0\longrightarrow M^q/J_n^q\longrightarrow H^q(X,F/I^nF)
\longrightarrow K_n^{q+1}\longrightarrow0.
\] Its compatible kernel coordinates vanish by the displayed transition bound; the remaining quotient filtration is cofinal with \(I^nM^q\). This proves the inverse-limit comparison, without an unexplained exchange of cohomology and inverse limit.

S02 Lemmas 1.2--1.3 prove exact finite-module completion, Noetherianity and faithful flatness of local completion. Thus a finite module over a complete local Noetherian ring is complete. S02 Lemma 1.4 proves unique lifting of idempotents through nilpotent ideals, and Corollary 2.4 turns the compatible thickened idempotents into an actual section of a projective scheme. These are precisely the two completion operations in the quotient proof below.

\phantomsection\label{AG-RG-03-ag-rg-03-grassmannian}

\textbf{Lemma 5.A (Grassmannians).} The scheme \(\operatorname{Gr}_d(V)\) of locally free rank-\(d\) quotients of a finite locally free \(V/S\) is smooth and projective, commutes with arbitrary base change, and its Plücker line bundle is \(\det Q\), for its universal quotient \(Q\).

\textbf{Proof.} Trivialize \(V=\mathcal O^n\). On the open where columns \(I\) give a basis of a quotient, its unique normalized matrix has those columns equal to the identity and \(d(n-d)\) free entries. A second basis \(J\) gives the transition \(M_J=(M_I|_J)^{-1}M_I\). These open affine charts represent the quotient functor on all tests and glue. The determinant quotient \(\bigwedge^dV\twoheadrightarrow\det Q\) defines the Plücker map to \(\mathbf P(\bigwedge^dV)\), with the quotient convention. Over the target chart \(p_I\ne0\), the ratios of minors containing a single replacement of an identity column recover, with their determinant signs, every free matrix entry. The corresponding map of coordinate rings is surjective. These target charts cover the projective bundle, so the Plücker map is a closed immersion, and its \(\mathcal O(1)\) pulls back to \(\det Q\). The construction glues for general \(V\) and uses only invertible minors, hence works over nonreduced bases. The subbundle convention follows by dualizing. \(\square\)

\subsection{Projective homogeneous quotients}\label{AG-RG-03-projective-homogeneous-quotients}

\textbf{Lemma 5.1.} Let \(H\) be a smooth closed subgroup with geometrically connected fibres of a smooth affine \(S\)-group \(K\) with geometrically connected fibres. Suppose \(N_K(H)=H\), and suppose every geometric quotient \(K_{\bar s}/H_{\bar s}\) is proper. Then the fppf quotient \(X=K/H\) is a smooth projective \(S\)-scheme, locally on \(S\). Its universal conjugate subgroup \(\mathcal H\subset K_X\) supplies a canonical relatively ample line bundle \(\det(\operatorname{Lie}\mathcal H)^*\).

\textbf{Proof.} Work first on an affine part \(S=\operatorname{Spec}R\). Descend \(K,H\) and their group laws and closed immersion to a finite-type Noetherian model \(S_0\), retaining smoothness and geometrically connected fibres as in the normalizer lemma. The normalizer is now of finite presentation and commutes with base change. Its equality with \(H\) descends after a further finite stage: equality of their finitely generated defining ideals uses finitely many relation witnesses. We do not assume that projectivity of all geometric quotient fibres descends merely by choosing coefficients.

Construct the quotient \(X_0=K_0/H_0\) on this model. The free right \(H_0\)-action defines a smooth equivalence relation. Corollary 8.3 of \href{https://kokunoyumeto.github.io/open-math-courses/courses/AG-RG/AG-RG-S07.html}{Flat quotient bootstrap} applies to this free smooth action: its quotient is an algebraic space, \(K_0\to X_0\) is an \(H_0\)-torsor, and its equality relation is \(H_0\times K_0=K_0\times_{X_0}K_0\) on every test scheme. Thus \(X_0\) is smooth with geometrically connected fibres, and formation commutes with base change. For the next two paragraphs write \(K,H,X\) for these model objects.

Let \(K_n,H_n\) be the \(n\)th infinitesimal neighbourhoods of the identities. The normalizer lemma identifies \(N_K(H)\) with the stabilizer of \(H_n\) in \(K_n\) for some \(n\), as functors on all test schemes. Increase \(n\) to at least one. Since \(N_K(H)=H\), this stabilizer is exactly \(H\). The coordinate modules of \(K_n\) and \(H_n\) are finite locally free, with augmentation graded pieces \(\operatorname{Sym}^i(\mathfrak k^*)\) and \(\operatorname{Sym}^i(\mathfrak h^*)\) for \(0\leq i\leq n\).

Conjugation acts on the Grassmannian of quotients of the locally free coordinate module of \(K_n\) of rank equal to that of \(H_n\). Its section corresponding to \(H_n\) has stabilizer exactly \(H\). Thus the orbit map gives a finite-type monomorphism \(X\to\operatorname{Gr}\). A finite-type monomorphism is quasi-finite and separated; the normalization lemma just proved from the earlier supporting lesson makes it quasi-affine. In particular \(X\) is a scheme and is quasi-projective locally on the base.

We verify properness near the image of the original \(S\) in \(S_0\). At any point \(p\) of this image, the quotient fibre is proper: after extension to an algebraic closure containing a residue field of a point above \(p\), this is the hypothesized proper geometric fibre, and properness descends under field extensions. Localize \(S_0\) at \(p\) and pass faithfully flatly to its completion \(A\). Place \(X\) as an open subscheme of a projective closure \(\bar X\) over \(A\). Its proper special fibre \(X_0\) is open and closed in \(\bar X_0\), and is therefore cut out by an idempotent. Idempotents lift uniquely across nilpotent ideals. The internal theorem on formal functions gives \(\Gamma(\bar X,\mathcal O_{\bar X})\simeq\varprojlim_n\Gamma(\bar X_n,\mathcal O_{\bar X_n})\), since \(A\) is complete and \(\bar X\) is proper. It therefore turns this compatible formal idempotent into an actual idempotent. It defines an open and closed proper component \(Z\subset\bar X\) with \(Z_0=X_0\). The closed subset \(Z\setminus X\) has empty special fibre; properness makes it empty. Hence \(Z\subset X\).

The image of \(Z\to\operatorname{Spec}A\) is both open and closed: it is proper, and it is the restriction of the smooth map \(X\to\operatorname{Spec}A\). It contains the closed point, so is all of the local base. In each geometric fibre, \(Z\) is a nonempty open and closed part of the connected fibre of \(X\). It is the entire fibre. Thus \(Z=X\), proving properness. This descends from the completion to \(\mathcal O_{S_0,p}\). The monomorphism to the Grassmannian is now proper there, hence a closed immersion. The finite-presentation limit theorem for closed immersions, proved in \href{https://kokunoyumeto.github.io/open-math-courses/courses/AG-MO/AG-MO-03.html}{Limits of schemes and Noetherian approximation}, Theorem 4.2, spreads this closed immersion to an open neighbourhood of \(p\). The union of these neighbourhoods contains the image of the original base. Pulling back to \(S\) proves the required projectivity there; it has not been imposed on unrelated model fibres.

The Grassmannian map is now a proper monomorphism and hence a closed immersion. Its Plücker bundle pulls back to \(\det\mathcal O_{\mathcal H_n}\). The augmentation filtration identifies this determinant with

\[
\bigotimes_{i=0}^n\det(\operatorname{Sym}^i(\operatorname{Lie}\mathcal H)^*)
=\bigl(\det(\operatorname{Lie}\mathcal H)^*\bigr)^c
\]

for a positive integer \(c\) when \(\dim H>0\). Positivity follows already from the \(i=1\) term; the determinant exponent for \(\operatorname{Sym}^i\) on a rank-\(d\) bundle is \(\binom{i+d-1}{d}\). Thus the stated canonical bundle is relatively ample. The zero-dimensional connected smooth case is trivial. All the constructions descend to the original arbitrary base and commute with base change. \(\square\)

The quotient proof is \href{https://kokunoyumeto.github.io/open-math-courses/courses/AG-RG/AG-RG-S07.html}{Flat quotient bootstrap}, Theorem 8.2 and Corollary 8.3; the formal-functions proof is \href{https://kokunoyumeto.github.io/open-math-courses/courses/AG-RG/AG-RG-S02.html}{Cohomology and models}, Theorem 2.3. The algebraic-space normalization lemma has its full earlier proof in the descent and Zariski Main supporting lesson, Theorem E9.1. The group-specific argument is original exposition.

\subsection{A universal finite cohomology complex}\label{AG-RG-03-a-universal-finite-cohomology-complex}

\textbf{Lemma (finite projective replacement).} Let \(A\) be Noetherian and let \(C\) be a complex of flat \(A\)-modules concentrated in degrees \([a,b]\), with finite cohomology modules. There is a complex \(K\) of finite projective modules in \([a,b]\) and a quasi-isomorphism \(K\to C\) that remains a quasi-isomorphism after tensoring with any \(A\)-module.

\textbf{Proof.} Construct a bounded-above complex \(P\) of finite free modules and a map \(f:P\to C\), descending from degree \(b\). Once the terms above \(j\) are chosen, its cone has terms

\[
D^i=C^i\oplus P^{i+1},\qquad
d(c,p)=(d_Cc+f(p),-d_Pp).
\]

Its cohomology is finite: the cone long exact sequence uses the finite cohomology of \(C\) and of the already chosen finite free portion of \(P\). Choose finitely many cocycles \((c_\ell,p_\ell)\) generating \(H^j(D)\). Take \(P^j\) free on \(e_\ell\), and set \(f(e_\ell)=c_\ell\), \(d_P(e_\ell)=-p_\ell\). The cocycle equations give \(d_P^2=0\) and \(d_Cf=fd_P\). In the new cone, \(d(0,e_\ell)=(c_\ell,p_\ell)\), so this kills \(H^j\) without changing higher cohomology. Continue to the left. Each degree stabilizes, giving a quasi-isomorphism \(P\to C\).

Its cone is acyclic, bounded above and flat termwise. Such a complex stays acyclic after every tensor: start at its last degree and descend through

\[
0\longrightarrow Z^i\longrightarrow D^i\longrightarrow Z^{i+1}\longrightarrow0.
\]

The quotient is flat; hence the kernel is flat and the sequence stays exact under tensor. Thus \(H^i(P\otimes_AM)=H^i(C\otimes_AM)\) for every \(M\).

Put \(Q=\operatorname{coker}(P^{a-1}\to P^a)\). The left tail is a free resolution of \(Q\), and

\[
\operatorname{Tor}_1^A(Q,M)=H^{a-1}(P\otimes_AM)=0
\]

for every \(M\), since \(C\otimes_AM\) starts in degree \(a\). The module \(Q\) is finitely presented, so it is flat and finite projective. Take \(K^a=Q\) and \(K^i=P^i\) for \(a<i\leq b\). The map to \(C\) factors through \(Q\) because \(C^{a-1}=0\). Truncation preserves cohomology, and the same acyclic flat-cone argument proves the assertion after every tensor. \(\square\)

Apply this to the finite section complex of a line bundle \(L_0\) on a smooth projective scheme over a Noetherian ring \(A_0\). Choose a finite affine cover. Separatedness makes its intersections affine, and S01 Sections 4--5 prove that its ordered section complex computes quasi-coherent cohomology. On each intersection the line bundle is a projective module over its chart algebra; smoothness makes that algebra \(A_0\)-flat, so every term is \(A_0\)-flat. S01 Theorem 8.4 and its closed-immersion passage prove that all cohomology modules are finite. Under every algebra base change, affine module-sheaf correspondence and localization, proved in S01 Section 2, identify the new section complex with its tensor product over \(A_0\). The finite projective replacement just proved therefore computes the line bundle cohomology after every base change.

\subsection{Sections in a vanishing family}\label{AG-RG-03-sections-in-a-vanishing-family}

\textbf{Lemma (sections in a vanishing family).} Let \(f:X\to S\) be smooth and projective of finite presentation, and let \(L\) be an invertible sheaf on \(X\). Suppose that \(H^q(X_{\bar s},L_{\bar s})=0\) for every geometric point \(\bar s\) and every \(q>0\). Then \(f_*L\) is finite locally free, all its positive higher direct images vanish, and these statements commute with every base change. In particular the canonical map \(g^*f_*L\to f'_*L'\) is an isomorphism for any \(g:S'\to S\). The base may be nonreduced.

\textbf{Proof.} Work with \(S=\operatorname{Spec}A\). Choose a projective embedding, a finite affine cover trivializing \(L\), and finite affine covers of the pairwise and triple overlaps. Write \(A\) as the filtered union of its finitely generated \(\mathbf Z\)-subalgebras. The embedding, chart opens and overlap maps descend at a finite stage by \href{https://kokunoyumeto.github.io/open-math-courses/courses/AG-MO/AG-MO-03.html\#4-descending-objects-maps-and-properties}{finite-presentation descent}, Theorems 4.1--4.2. Retain the closed immersion, so the model \(X_0\) is projective.

Retain smoothness at a finite stage by S02 Lemma 3.1 and Proposition 3.2. In a finite polynomial presentation, Lemma 3.0 identifies formal smoothness with a left inverse of the conormal differential map. The left inverse uses finitely many coefficients; its equations on the finitely many relation generators have finite polynomial witnesses. Descend these coefficients and witnesses to one stage. The same left inverse then proves formal smoothness there, and the retained finite presentation proves smoothness. Apply this on a finite affine cover, retaining the cover and the projective closed immersion by the finite-presentation limit proof. The line-bundle transition units, their inverses and the finitely many overlap and cocycle witnesses descend at the same stage, as in S02 Proposition 3.2. Gluing these descended transition units defines the invertible sheaf \(L_0\) on the smooth projective model \(X_0/A_0\).

The preceding finite-projective-replacement lemma gives a bounded finite projective complex \(K_0\) such that

\[
H^q(K_0\otimes_{A_0}B)
\simeq H^q(X_0\times_{A_0}\operatorname{Spec}B,L_0\otimes_{A_0}B)
\]

for every \(A_0\)-algebra \(B\), compatibly with the canonical comparison maps. Its proof uses the finite affine-cover complex and projective Serre finiteness; ordinary tensor of \(K_0\) computes its derived tensor. This finite argument supplies the universal comparison before pullback to the original base.

Put \(K=K_0\otimes_{A_0}A\). Near any \(s\in\operatorname{Spec}A\), trivialize its terms. Whenever a differential has an entry invertible at \(s\), shrink the neighbourhood and perform row and column operations to put that entry in an identity block. The equation \(d^2=0\) makes the adjacent differentials zero on that block, so it splits off a contractible pair \(A\xrightarrow{1}A\). After finitely many cancellations, all remaining differential entries vanish in \(\kappa(s)\). The remaining complex therefore has

\[
K^q\otimes_A\kappa(s)=H^q(X_s,L_s).
\]

For \(q>0\) the right side is zero by the hypothesis and field extension; for \(q<0\) it is zero because \(L_s\) is a sheaf. Thus every remaining free term except degree zero has rank zero and disappears on a neighbourhood of \(s\). The degree-zero term is finite free. Every cancelled identity pair stays contractible after any tensor. The universal comparison consequently proves the direct-image assertions and every base-change assertion, including nonflat changes involving nilpotents. The canonical comparisons glue. Notice that vanishing was imposed only on the original family, not on unrelated fibres of its Noetherian model. \(\square\)

\textbf{Proposition 5.2.} If \(G\) has trivial centre and geometric semisimple rank one, and has a split maximal torus \(T\), then Zariski locally \((G,T)\) is isomorphic to \((\operatorname{PGL}_2,D)\).

\textbf{Proof.} Its geometric fibres are \(\operatorname{PGL}_2\) by Proposition 2.3: the action on the projective line has trivial kernel. Choose a root \(\alpha\) locally. It is an isomorphism \(T\to\mathbf G_m\), as this is true on the character lattices of the geometric fibres. Put \(B=P_G(\alpha^{-1})\). It is smooth and on each geometric fibre is an upper triangular Borel subgroup of \(\operatorname{PGL}_2\).

Its scheme-theoretic normalizer equals \(B\). The normalizer lemma makes it closed of finite presentation and proves that its formation commutes with arbitrary base change. On geometric fibres the point stabilizer in \(\mathbf P^1\) is self-normalizing. Infinitesimally, the bracket of a diagonal matrix class with the lower off-diagonal class is a nonzero lower off-diagonal class in every characteristic, so the normalizer has exactly the tangent space of \(B\). It is therefore smooth and equals \(B\) on those fibres. The flat-source criterion above, applied with source \(B\), proves the equality over the base without assuming the normalizer flat.

Lemma 5.1 makes \(f:X=G/B\to S\) a smooth projective curve whose geometric fibres are \(\mathbf P^1\), and it has the section \(e=1B\). A section of a smooth relative curve is an effective Cartier divisor: étale coordinates at the section give one parameter cutting out its ideal, a nonzerodivisor because the coordinate algebra is flat over the polynomial chart. Put \(L=\mathcal O_X(e)\); it has degree one on every geometric fibre. The sections-in-a-vanishing-family lemma applies to \(L\), since \(H^1(\mathbf P^1,\mathcal O(1))=0\), and gives \(E=f_*L\) locally free of rank two, with arbitrary base change. The evaluation \(f^*E\to L\) is surjective: its cokernel is finite locally, its geometric fibres vanish by generation of \(\mathcal O(1)\), and local Nakayama kills that cokernel. With the quotient convention it defines \(j:X\to\mathbf P(E)\), and \(j_{\bar s}\) is the usual isomorphism given by the two sections of \(\mathcal O(1)\). Both schemes are smooth over \(S\), and the differential of \(j\) is an isomorphism on all geometric fibres, hence locally an isomorphism of their cotangent bundles. The étale differential criterion makes \(j\) étale. Its geometric fibres are isomorphisms, so it is universally injective and surjective; an étale universally injective morphism is an open immersion, and surjectivity makes it an isomorphism. Trivializing \(E\) Zariski locally now gives \(G/B\simeq\mathbf P^1_S\).

The action now gives \(G\to\operatorname{PGL}_2\), an isomorphism on every geometric fibre, hence an isomorphism over \(S\). It carries \(B\) to the point stabilizer. Within that stabilizer the transporter taking its torus \(T\) to the diagonal torus is a \(\mathbf G_m\)-torsor: the stabilizer of the diagonal torus inside the upper triangular group is the diagonal torus. Such a torsor is the frame bundle of a line bundle and is Zariski locally trivial. This finishes the pairwise assertion. \(\square\)

\section{6. Splitting a central extension of the rank-one cover}\label{AG-RG-03-splitting-a-central-extension-of-the-rank-one-cover}

We need an integral result, including characteristic two. The following splitting argument is due to Gabber; we give the argument for a separated commutative group scheme \(Z\), which includes the multiplicative-type centres used here.

\textbf{Theorem 6.1.} Every fppf central extension

\[
1\longrightarrow Z\longrightarrow E\longrightarrow\operatorname{SL}_{2,S}\longrightarrow1
\]

has a unique homomorphic section.

\textbf{Proof.} In \(\operatorname{SL}_2\) use

\[
x(u)=\begin{pmatrix}1&u\\0&1\end{pmatrix},\quad
y(v)=\begin{pmatrix}1&0\\v&1\end{pmatrix},\quad
h(t)=\operatorname{diag}(t,t^{-1}).
\]

Let \(D',U',V'\) be their inverse images in \(E\). A central extension of a commutative group has a commutator pairing on its quotient: changing a lift by an element of the central kernel does not change the commutator, and the identities for commutators make this pairing additive in each variable.

For \(D'\), the pairing \(\mathbf G_m\times\mathbf G_m\to Z\) is trivial. Restrict the first argument to \(\mu_n\). The resulting pairing is killed by \(n\) in that argument, and the \(n\)th-power map is an fppf epimorphism in the second argument, so this restriction is trivial. The schemes \(\mu_n\) are universally schematically dense in \(\mathbf G_m\): distinct Laurent exponents remain distinct modulo sufficiently large \(n\). Separatedness of \(Z\) makes the whole pairing trivial. Thus \(D'\) is commutative.

The conjugation action on \(E\) factors through \(\operatorname{SL}_2\), since \(Z\) is central. The commutator pairing on \(U'\) is invariant under simultaneous multiplication of its arguments by \(t^2\). The square map is fppf surjective, so it is invariant under simultaneous multiplication by every unit. For given \(u,v\), the morphism \(a\mapsto c(au,av)\) from \(\mathbf A^1\) to \(Z\) is constant on \(\mathbf G_m\). Separatedness extends the equality to zero, where it is the identity. Hence \(c(u,v)=1\). Thus \(U'\) is commutative; the same argument proves it for \(V'\).

We construct a unique \(D\)-equivariant section of \(U'\to U\). Fppf locally choose a unit \(t\) with \(t^2-1\) invertible and a lift \(\tilde h\) of \(h(t)\). Such choices exist: the open locus \(t(t^2-1)\ne0\) in an affine line is smooth surjective, and the extension is an fppf epimorphism. In the commutative group \(U'\), the endomorphism

\[
\varphi(a)=\tilde h a\tilde h^{-1}a^{-1}
\]

kills \(Z\) and induces multiplication by \(t^2-1\) on \(U=\mathbf G_a\). It therefore factors through a homomorphism \(U\to U'\), which after division by \(t^2-1\) is the required section. It is \(D\)-equivariant because \(D'\) is commutative. Uniqueness holds since any difference is a \(D\)-invariant homomorphism \(U\to Z\): scalar invariance and the specialization from nonzero scalars to zero force it to vanish. Uniqueness descends these local sections. Obtain \(x':\mathbf G_a\to E\) and similarly \(y'\).

Set

\[
w'(t)=y'(-t^{-1})x'(t)y'(-t^{-1}),\qquad
h'(t)=w'(t)w'(1)^{-1}.
\]

Their images are \(\begin{pmatrix}0&t\\-t^{-1}&0\end{pmatrix}\) and \(h(t)\). Thus \(h'(t)\in D'\). Equivariance gives

\[
h'(s)w'(t)h'(s)^{-1}=w'(s^2t).
\]

Divide by the instance \(t=1\). Since \(D'\) is commutative, this gives \(h'(t)=h'(s^2t)h'(s^2)^{-1}\). The square map being fppf surjective, \(h'\) is a homomorphism. Its definition then gives \(h'(s)w'(t)=w'(st)\).

Substitute the expression for \(w'\) into this last equality and move the \(y'\) terms using \(h'(s)y'(v)h'(s)^{-1}=y'(s^{-2}v)\). With \(u=st\) and \(v=(s-1)/(st)\), the result is

\[
x'(u)y'(v)=y'\left(\frac{v}{1+uv}\right)
h'(1+uv)x'\left(\frac{u}{1+uv}\right)
\tag{1}
\]

when \(u\) and \(1+uv\) are units. The restriction that \(u\) be a unit can be removed. Fppf locally write \(u=a+b\) with \(a,b,1+bv\) units; the open conditions on a choice of \(b\) omit only finitely many points of each geometric affine-line fibre. Apply (1) first to \(x'(b)y'(v)\), then to \(x'(a)y'(v/(1+bv))\). The second denominator is \((1+uv)/(1+bv)\). Moving the two resulting \(h'\) factors together leaves the \(y'\) parameter \(v/(1+uv)\) and the \(x'\) parameter

\[
\frac{a}{(1+bv)(1+uv)}+\frac{b}{1+bv}
=\frac{u}{1+uv}.
\]

This proves (1) in its full domain.

These identities, the additivity of \(x',y'\), the multiplicativity of \(h'\), and its two conjugation identities are exactly the multiplication rules on the open cell \(\Omega=U_-DU_+\) of \(\operatorname{SL}_2\). Define a section there by

\[
y(v)h(t)x(u)\longmapsto y'(v)h'(t)x'(u).
\]

It is multiplicative whenever both factors and their product are in \(\Omega\): reorder their factors with (1). Write \(s_\Omega:\Omega\to E\) for this cell section. For any test scheme \(R\to S\) and finite list \(g_1,\ldots,g_m\in\operatorname{SL}_2(R)\), the open \[
\Omega_R\cap\bigcap_i\Omega_Rg_i^{-1}
\] is nonempty in every geometric fibre: a smooth geometrically connected group has integral geometric fibres, and finitely many dense opens have nonempty intersection. It is smooth and surjective over \(R\); affine refinements therefore supply an fppf local point \(c\) with \(c,cg_i\in\Omega\).

In particular, for one \(g\), use the nonempty open \[
\Omega_R\cap\Omega_R^{-1}\cap\Omega_Rg^{-1}.
\] Then \(g=c^{-1}(cg)\) is a product of two points of \(\Omega\). Assign to a word in cell points the corresponding product of their \(s_\Omega\)-values. This is independent of the word. For two words with the same endpoint \(g\), choose \(c\) so that \(c\) and its translates by every partial product of both words lie in \(\Omega\). Successive cell multiplicativity says that multiplying either lifted word on the left by \(s_\Omega(c)\) gives \(s_\Omega(cg)\). Cancellation proves equality. In particular a word representing the identity has lift one; concatenation therefore respects multiplication and inverses. These word values agree after further fppf pullbacks, so descend through the sheaf property of \(E\). They define a homomorphism of represented fppf sheaves, hence a scheme homomorphism \(\operatorname{SL}_2\to E\). Its composite with \(E\to\operatorname{SL}_2\) is the identity on generating words. It is the required homomorphic section on all test schemes.

Finally any homomorphism \(\operatorname{SL}_2\to Z\) kills \(x(u)\), since

\[
[h(t),x(u)]=x((t^2-1)u),
\]

and one may choose \(t^2-1\) invertible fppf locally. It kills \(y(v)\) similarly, and then \(h(t)=w(t)w(1)^{-1}\). It kills the open cell, hence the group. The difference between two central sections is such a homomorphism, proving uniqueness. \(\square\)

\section{7. Coroots from multiplication}\label{AG-RG-03-coroots-from-multiplication}

\textbf{Theorem 7.1.} For every root \(\alpha\) there is a canonical coroot \(\alpha^\vee:\mathbf G_m\to T\), with \(\langle\alpha,\alpha^\vee\rangle=2\). There is a canonical perfect pairing \(\mathfrak g_\alpha\otimes\mathfrak g_{-\alpha}\to\mathcal O_S\) such that, writing its value as \(XY\),

\[
\exp_\alpha(X)\exp_{-\alpha}(Y)
=\exp_{-\alpha}\left(\frac{Y}{1+XY}\right)
\alpha^\vee(1+XY)
\exp_\alpha\left(\frac{X}{1+XY}\right)
\tag{2}
\]

whenever \(1+XY\) is a unit. This unit condition is precisely membership in the open cell \(U_{-\alpha}TU_\alpha\).

\textbf{Proof.} Work in \(G_\alpha\), whose scheme-theoretic centre is \(Z=\ker\alpha\subset T\), by Section 1. Theorem 1.1 constructs \(q:G_\alpha\to Q=G_\alpha/Z\) as an affine reductive quotient and an fppf \(Z\)-torsor, compatible with every base change. The torus quotient has character group \(\mathbf Z\alpha\subset X^*(T)\), so \(T/Z\simeq\mathbf G_m\) via the descended character \(\alpha\). It is a closed subgroup of \(Q\): its pullback under \(q\) is \(T\subset G_\alpha\), and closed immersions descend through the fppf torsor. On geometric fibres Proposition 2.3 identifies this quotient with \(\operatorname{PGL}_2\), so \(T/Z\) is a maximal torus.

The open cell \(U_{-\alpha}TU_\alpha\) is right \(Z\)-stable; centrality makes \(Z\) act only on its torus coordinate. Its quotient is the open subscheme \[
U_{-\alpha}\times(T/Z)\times U_\alpha\subset Q,
\] as follows by invariant-open descent and the product torsor. At the identity this gives the two nonzero torus weights \(\alpha,-\alpha\), with their root lines, in \(\operatorname{Lie}Q\). The centre formula of Section 1, applied to \(Q,T/Z\), is consequently \(Z(Q)=\ker(\alpha:T/Z\to\mathbf G_m)=1\), as a scheme over the original base. Proposition 5.2 therefore identifies the pair \(Q,T/Z\), Zariski locally, with \(\operatorname{PGL}_2\) and its diagonal torus. Choose the identification so \(\alpha(\operatorname{diag}(t,1))=t\).

Pull back \(G_\alpha\to\operatorname{PGL}_2\) by \(\operatorname{SL}_2\to\operatorname{PGL}_2\). Theorem 6.1 splits the resulting central extension by \(Z\). The homomorphism \(\operatorname{SL}_2\to G_\alpha\) obtained from the splitting identifies its upper and lower unipotent groups with \(U_\alpha,U_{-\alpha}\): the maps are equivariant, and their differentials on the root lines are isomorphisms, so Theorem 4.1 gives the assertion. The diagonal cocharacter therefore defines \(\alpha^\vee\).

Direct multiplication gives

\[
\begin{pmatrix}1&x\\0&1\end{pmatrix}
\begin{pmatrix}1&0\\y&1\end{pmatrix}
=
\begin{pmatrix}1&0\\y/(1+xy)&1\end{pmatrix}
\begin{pmatrix}1+xy&0\\0&(1+xy)^{-1}\end{pmatrix}
\begin{pmatrix}1&x/(1+xy)\\0&1\end{pmatrix}.
\]

The upper-left entry is a unit exactly when the product is in this open cell. It supplies (2), a perfect pairing of the two root lines, and \(\alpha\circ\alpha^\vee(t)=t^2\).

The pairing and coroot are unique. In the adjoint quotient, comparing the upper-root parameter in (2) forces any alternative scalar in the pairing to be one; equality of open-cell coordinates then fixes the coroot on every unit \(1+xy\), which ranges over all units by setting \(y=1\). An alternative coroot in \(T\) with the same adjoint image could differ by a cocharacter into \(Z\), but (2) fixes its actual torus component and eliminates that difference. Hence the constructions on different local identifications agree and descend. The same characterization proves base-change compatibility. \(\square\)

For \(\operatorname{SL}_2\), use \(T\simeq\mathbf G_m\) via \(\operatorname{diag}(t,t^{-1})\). Its root is \(\alpha(t)=t^2\), and \(\alpha^\vee(t)=\operatorname{diag}(t,t^{-1})\). For \(\operatorname{PGL}_2\), parameterize its diagonal torus by the class of \(\operatorname{diag}(t,1)\). Its root is \(t\mapsto t\), and its coroot is the class of \(\operatorname{diag}(t^2,1)\). In both cases the pairing is two, but the root and coroot occupy different sublattices.

For \(\operatorname{GL}_2\), the roots are \(e_1-e_2\) and its negative, and the positive coroot is \(t\mapsto\operatorname{diag}(t,t^{-1})\). These explicit formulas remain valid over \(\mathbb Z\).

\section{8. The three rank-one groups}\label{AG-RG-03-the-three-rank-one-groups}

\subsection{A lattice calculation with pairing two}\label{AG-RG-03-a-lattice-calculation-with-pairing-two}

\textbf{Lemma 8.A.} Let \(X\) be a finite free abelian group of rank \(r+1\). Let \(a\in X\) and \(b\in X^*=\operatorname{Hom}(X,\mathbf Z)\) be nonzero and satisfy \(b(a)=2\). Up to an isomorphism of \(X\) preserving this pair, exactly one of the following descriptions occurs:

\begin{enumerate}
\def\labelenumi{\arabic{enumi}.}
\tightlist
\item
  \(X=\mathbf Z e_0\oplus\mathbf Z^r\), \(a=2e_0\), \(b=e_0^*\).
\item
  \(X=\mathbf Z e_0\oplus\mathbf Z^r\), \(a=e_0\), \(b=2e_0^*\).
\item
  \(r\geq1\), \(X=\mathbf Z e_1\oplus\mathbf Z e_2\oplus\mathbf Z^{r-1}\), \(a=e_1-e_2\), \(b=e_1^*-e_2^*\).
\end{enumerate}

In these formulas \(b\) is zero on the displayed extra summand.

\textbf{Proof.} Write \(a=m a_0\) and \(b=n b_0\), with \(a_0,b_0\) primitive and \(m,n>0\). Then

\[
mn\,b_0(a_0)=2.
\]

Consequently \((m,n)\) is \((2,1)\), \((1,2)\), or \((1,1)\).

In the first case \(b(a_0)=1\). Every \(x\in X\) has the unique expression

\[
x=b(x)a_0+\bigl(x-b(x)a_0\bigr),
\qquad x-b(x)a_0\in\ker b.
\]

Thus \(X=\mathbf Z a_0\oplus\ker b\); choose a basis of the free kernel. This gives case 1. In the second case apply the same expression using \(b_0(a)=1\). It gives case 2.

Suppose now that both \(a\) and \(b\) are primitive. Extend \(a\) to a basis \(a,c_1,\ldots,c_r\). At least one \(b(c_i)\) is odd: otherwise every value of \(b\), including \(b(a)=2\), would be even, contradicting primitivity of \(b\). Relabel so that \(b(c_1)=2q+1\), and replace \(c_1\) by \(c_1-qa\). This is a basis operation and makes its value under \(b\) equal to one. Call the new vector \(c\). Replace each remaining \(c_i\) by \(c_i-b(c_i)c\); their values become zero. Finally set

\[
e_1=c,\qquad e_2=c-a.
\]

These replace the basis pair \(a,c\) by another unimodular basis pair. They satisfy \(a=e_1-e_2\), \(b(e_1)=1\) and \(b(e_2)=-1\), giving case 3. If \(r=0\), both primitive elements would pair to \(\pm1\), so this case cannot occur.

The divisibilities \((m,n)\) of \(a\) and \(b\) are invariant under lattice isomorphisms. They distinguish all three cases. \(\square\)

\subsection{Extending a map from the rank-one cell}\label{AG-RG-03-extending-a-map-from-the-rank-one-cell}

\textbf{Lemma 8.B.} Let \(K/S\) be a smooth group with geometrically connected fibres. Let \(\Omega\subset K\) be a fibrewise dense open containing the identity. Let \(H/S\) be a group scheme. Suppose \(f:\Omega\to H\) sends the identity to the identity and satisfies

\[
f(xy)=f(x)f(y)
\]

on the open where \(x,y,xy\in\Omega\). Then \(f\) extends uniquely to a homomorphism \(K\to H\).

\textbf{Proof.} A geometrically connected smooth algebraic group is geometrically irreducible: its regular irreducible components are disjoint and open, so connectedness leaves only one. For any test scheme \(R\to S\) and any finite list \(g_1,\ldots,g_d\in K(R)\), the open

\[
\Omega_R\cap\bigcap_{i=1}^d\Omega_R g_i^{-1}
\]

has a nonempty open in every geometric fibre of \(K_R\). Its map to \(R\) is smooth and surjective. It therefore supplies a point \(c\) after an fppf cover of \(R\), with \(c\) and all \(cg_i\) in the cell.

In particular, for \(g\in K(R)\), choose \(c\in\Omega_R\cap\Omega_R^{-1}\cap\Omega_Rg^{-1}\) locally. Then \(g=c^{-1}(cg)\) is a product of two cell points. Assign to any word in cell points the corresponding product of \(f\)-values.

This assignment is independent of the word. For two words \(x_1\cdots x_d\) and \(y_1\cdots y_e\) with the same endpoint \(g\), choose \(c\) so that it and every translated partial product of both words lie in the cell. Repeatedly apply the given identity. The first image word multiplied on the left by \(f(c)\), and the second image word multiplied by the same element, both equal \(f(cg)\). Cancellation proves equality. The assignment also respects inverses: for \(x\in\Omega(R)\), choose \(d,d^{-1},dx^{-1}\in\Omega\). The identities \((dx^{-1})x=d\) and \(d^{-1}d=1\) then give \(f(d^{-1})f(dx^{-1})=f(x)^{-1}\). The word \(d^{-1}(dx^{-1})\) represents \(x^{-1}\), as required.

The assignments agree after fppf pullback, so the sheaf property of \(H\) descends them. Concatenation of words proves multiplicativity. Yoneda makes this map of represented fppf sheaves a scheme morphism. Every extension must have these values on generating words, proving uniqueness. \(\square\)

\subsection{The three rank-one models}\label{AG-RG-03-the-three-rank-one-models}

\textbf{Theorem 8.1.} A split reductive group of rank \(r+1\) and semisimple rank one over a field is isomorphic to exactly one of

\[
\mathbf G_m^r\times\operatorname{SL}_2,
\qquad
\mathbf G_m^r\times\operatorname{PGL}_2,
\qquad
\mathbf G_m^{r-1}\times\operatorname{GL}_2\quad(r\geq1).
\]

For a reductive group over a scheme with a split maximal torus and semisimple rank one, the same description holds Zariski locally on fixed-type loci. Trivial root lines with chosen frames give a global isomorphism to the corresponding base change from \(\mathbf Z\).

\textbf{Proof.} The root and coroot construction in Theorem 7.1 gives two roots \(\pm\alpha\), their additive root lines, the coroot \(\eta=\alpha^\vee\), and \(\langle\alpha,\eta\rangle=2\). Apply Lemma 8.A to \(X=X^*(T)\), \(a=\alpha\), \(b=\eta\). The three resulting pairs are precisely those computed for the three displayed matrix groups: in \(\operatorname{SL}_2\) the root is twice the primitive character and the coroot is primitive; in \(\operatorname{PGL}_2\) the root is primitive and the coroot twice primitive; in \(\operatorname{GL}_2\) they are respectively \(e_1-e_2\) and \(e_1^*-e_2^*\).

Let \(G_0\) be the indicated matrix model. The lattice isomorphism gives a torus isomorphism \(f_T:T\to T_0\) carrying root characters and coroots to those of the model. Choose a nonzero positive-root frame \(E\in\mathfrak g_\alpha\). The perfect pairing of Theorem 7.1 gives a unique negative-root frame \(F\) with pairing one. Define \(x(u)=\exp_\alpha(uE)\) and \(y(v)=\exp_{-\alpha}(vF)\), and match these to the upper and lower matrix parameters \(x_0,y_0\) of \(G_0\).

The open cell theorem gives \(\Omega=U_-TU_+\). On it define

\[
f\bigl(y(v)t x(u)\bigr)=y_0(v)f_T(t)x_0(u).
\]

This is an isomorphism of schemes from \(\Omega\) to the model cell. The group law on the torus, the additivity of \(x,y\), and the conjugation identities \(t x(u)t^{-1}=x(\alpha(t)u)\), \(t y(v)t^{-1}=y(\alpha(t)^{-1}v)\), agree under this map. The remaining multiplication rule is

\[
x(u)y(v)=y\left(\frac{v}{1+uv}\right)
\eta(1+uv)x\left(\frac{u}{1+uv}\right),
\qquad 1+uv\in\mathcal O^\times,
\]

proved in Theorem 7.1 by the arbitrary-base central-extension argument and the explicit \(2\times2\) matrix identity.

For two cell points \(y(v)t x(u)\) and \(y(v')t'x(u')\), their product is in the cell exactly when \(x(u)y(v')\) is: left multiplication by \(U_-T\) and right multiplication by \(TU_+\) preserve the cell, as do their inverses. The preceding rule and its unit criterion therefore give every multiplication of cell points whose product remains in the cell. It follows that \(f\) satisfies Lemma 8.B. The lemma extends it to a group homomorphism \(G\to G_0\). The inverse cell map satisfies the same identities and extends in the other direction. Uniqueness in the lemma makes their compositions the identity. Thus \(G\simeq G_0\).

Any group isomorphism carries a maximal torus to a maximal torus. Over an algebraic closure the torus conjugacy theorem then identifies its root/coroot pair with the pair for the fixed torus. It can interchange the signs of both elements, but does not change either divisibility. The three invariant pairs \((2,1),(1,2),(1,1)\) distinguish the models in all characteristics, including characteristic two. This proves uniqueness without a forward reference to the derived group.

Over a scheme, restrict to an open where the root/coroot lattice pair is constant and the positive root line is trivial. Its paired negative line is then also trivial. The same torus map, cell formula and fppf word argument work over every test scheme, so give the asserted local isomorphism. When the lattice pair and frames are fixed globally, these constructions are global. All maps commute with base change because the displayed formulas and their character pairings do. \(\square\)

\section{9. Exercises and solutions}\label{AG-RG-03-exercises-and-solutions}

\textbf{Exercise 9.1 (first steps).} Compute all roots and root spaces of \(\operatorname{GL}_n\) with respect to its diagonal torus. Determine the zero-weight space and the positive coroot for \(e_i-e_j\).

\textbf{Solution.} The diagonal matrix units have weight zero and form \(\mathfrak t\). For \(i\ne j\), \(E_{ij}\) has weight \(e_i-e_j\), each once. The corresponding root group is \(I+uE_{ij}\). Its coroot has diagonal entries \(t\) in position \(i\), \(t^{-1}\) in position \(j\), and \(1\) elsewhere. Evaluating \(e_i-e_j\) gives \(t^2\), so the root--coroot pairing is two.

\textbf{Exercise 9.2 (rank-one multiplication).} Compute \(U_{\pm\alpha}\) and \(\alpha^\vee\) for \(\operatorname{SL}_2\). Verify the factorization (2), and explain why the character differential in characteristic two does not change this factorization.

\textbf{Solution.} The two groups are the upper and lower unipotent matrices. Their parameters are \(x,y\) in the displayed matrix identity, and \(\alpha^\vee(t)=\operatorname{diag}(t,t^{-1})\). Multiplication of the three factors gives \(\begin{pmatrix}1+xy&x\\y&1\end{pmatrix}\), equal to \(x_\alpha(x)x_{-\alpha}(y)\). The character on the upper group is \(t^2\), so its pairing with this coroot is two. Although its differential is zero in characteristic two, the Laurent character and all polynomial matrix identities remain valid; no division by two was used.

\textbf{Exercise 9.3 (infinitesimal information).} In characteristic different from two, show that \(\mathfrak{sl}_2\simeq\mathfrak{pgl}_2\) as Lie algebras, while the groups are not isomorphic. What happens to the differential of \(\operatorname{SL}_2\to\operatorname{PGL}_2\) in characteristic two?

\textbf{Solution.} The inclusion of trace-zero matrices followed by quotienting by scalar matrices has zero kernel when two is invertible: a trace-zero scalar is zero. Both Lie algebras have dimension three, so this bracket-preserving map is an isomorphism. The groups have different centres, \(\mu_2\) and \(1\). In characteristic two the scalar identity has trace zero and lies in the kernel of the differential, which therefore has one-dimensional kernel. It cannot be an isomorphism. The group map remains a central fppf isogeny with kernel \(\mu_2\).

\textbf{Exercise 9.4 (classification).} List the rank-two split reductive groups of semisimple rank one. Distinguish them by their centres and derived groups, including characteristic two.

\textbf{Solution.} Put \(r=1\) in Theorem 8.1. The list is \(\mathbf G_m\times\operatorname{SL}_2\), \(\mathbf G_m\times\operatorname{PGL}_2\), and \(\operatorname{GL}_2\). Their centres are respectively \(\mathbf G_m\times\mu_2\), \(\mathbf G_m\), and \(\mathbf G_m\). The first centre is not a torus: in characteristic two it is nonsmooth, and otherwise it has two geometric components. We compute the indicated derived groups directly. Fppf locally choose \(t\) with \(t^2-1\) invertible. The identities \([h(t),x(u)]=x((t^2-1)u)\) and \([h(t),y(v)]=y((t^{-2}-1)v)\) put both root groups of \(\operatorname{SL}_2\) in its commutator subgroup. The open cell word argument of Theorem 6.1 shows that these root groups generate \(\operatorname{SL}_2\) as an fppf sheaf, since \(h(t)=w(t)w(1)^{-1}\). Their images generate \(\operatorname{PGL}_2\), because \(\operatorname{SL}_2\to\operatorname{PGL}_2\) is fppf surjective; hence its derived group is itself. For \(\operatorname{GL}_2\), determinant kills all commutators, and conjugation by \(\operatorname{diag}(t,1)\), with \(t-1\) invertible fppf locally, puts both determinant-one root groups in the commutator subgroup. The same generation argument gives its derived group \(\operatorname{SL}_2\). The central \(\mathbf G_m\) factors contribute no commutators. Thus the last two models have derived groups \(\operatorname{PGL}_2\) and \(\operatorname{SL}_2\); the first and third have the same derived group but different centres. These computations are scheme-theoretic and remain valid in characteristic two.

\textbf{Exercise 9.5 (families).} For \(G=\operatorname{SL}(\mathcal O\oplus L)\), identify the two root groups and their pairing. Give a condition on \(L\) necessary for a pinning of this pair with its displayed torus.

\textbf{Solution.} An upper off-diagonal entry is a map \(L\to\mathcal O\), so its line is \(L^{-1}\); a lower entry is a map \(\mathcal O\to L\), so its line is \(L\). Composition gives \(L^{-1}\otimes L\to\mathcal O\), the perfect pairing in Theorem 7.1. A pinning chooses a nowhere-vanishing section in the simple root line. It therefore requires \(L^{-1}\), equivalently \(L\), to be trivial. This illustrates the global hypothesis absent from the field case.

The \href{https://kokunoyumeto.github.io/open-math-courses/courses/AG-RG/prerequisites.html}{course prerequisite guide} records the exact supporting statements, their full proof routes, and the hypotheses needed in their applications.

\section{References and prerequisite proofs}\label{AG-RG-03-references-and-prerequisite-proofs}

\begin{itemize}
\tightlist
\item
  Brian Conrad, \href{https://math.stanford.edu/~conrad/papers/luminysga3.pdf}{\emph{Reductive group schemes}}, §§2.3 and 4.1--4.3, for the dynamic approach and Gabber's central-extension splitting argument. The proofs above are written out independently.
\item
  \href{https://www.jmilne.org/math/CourseNotes/iAG200.pdf}{Milne's freely accessible \emph{Algebraic Groups}, version 2.00 (2015)}. This exact free author edition provides comparison material; its citations do not replace the programme proofs.
\item
  SGA 3, \href{https://webusers.imj-prg.fr/~patrick.polo/SGA3/}{\emph{Schémas en groupes}}, Exposés XIX--XX, for roots and the relative rank-one theorem.
\item
  The \href{https://stacks.math.columbia.edu/}{official Stacks project} and AI Integrated Stacks Project, an edition with AI-proposed corrections and AI-written additions, not reviewed by the Stacks project's maintainers. The normalization lemma in Section 5 and its proof follow the Stacks project authors' argument (Tag 0ABS) in this edition.
\end{itemize}

The following supporting statements have their exact proof routes recorded in the course prerequisite guide:

\begin{itemize}
\tightlist
\item
  \href{https://kokunoyumeto.github.io/open-math-courses/courses/AG-RG/AG-RG-S01.html}{Algebra and sheaf cohomology before reductive groups}, §§3--8: injectives, finite-cover comparisons, affine vanishing, projective-space twists, ample embeddings and projective Serre finiteness.
\item
  \href{https://kokunoyumeto.github.io/open-math-courses/courses/AG-RG/AG-RG-S02.html}{Projective cohomology and smooth affine models}, §§1--6: exact and faithful local completion; projective formal functions and idempotents; smooth connected affine models; universal finite projective complexes and all-base-change direct images; regular curve immersions, curve duality, divisor degrees and the genus-zero calculation.
\item
  \href{https://kokunoyumeto.github.io/open-math-courses/courses/AG-RG/AG-RG-S03.html}{Completing a smooth affine curve and extending its group action}, §§1--4: finite curve normalization, completion, the actual family action and one-dimensional unipotent identification used through the preceding solvable-group proofs.
\item
  \href{https://kokunoyumeto.github.io/open-math-courses/courses/AG-RG/AG-RG-S04.html}{Affine descent, Zariski Main and recognition of spaces}, Theorems E8.1--E9.1: recognition on every scheme test, representability and the full relative integral-closure factorization.
\item
  \href{https://kokunoyumeto.github.io/open-math-courses/courses/AG-RG/AG-RG-S07.html}{Flat quotient bootstrap}, Theorem 8.2 and Corollary 8.3: the algebraic-space quotient and torsor relation for the free smooth subgroup action, over arbitrary bases.
\item
  \hyperref[AG-RG-03-ag-rg-03-grassmannian]{Lemma 5.A of this lesson}: quotient Grassmannian charts, arbitrary base change, and the closed Plücker immersion with determinant quotient bundle; the subbundle convention follows by dualizing.
\item
  The current written group-scheme proofs: \emph{Diagonalizable groups}, Theorems 3.4--3.5 and §4, for smooth fixed loci, invariant algebras and torus quotients; \emph{Group schemes over a field}, Proposition 5.8, for the reduced schematic image and faithful-flat image factorization; \emph{Lie algebras and smoothness}, Theorem 2.9 and Theorem 3.2, for the universal normalizer tangent calculation and geometric smoothness criterion; and \emph{Quotients and torsors}, its affine descent proof and Theorem 8.1, for effective algebra descent and line-bundle frame torsors. These are the actual written proofs, including positive-characteristic and nonreduced scope, rather than source citations.
\end{itemize}

Theorem 1.1 proves the affine central quotient before the rank-one argument; the following lesson adds the consequences for root-group coordinates. The arbitrary-base derived group is constructed after the existence theorem in the following lesson.

\section{History}\label{AG-RG-03-history}

This lesson's exposition, constructions, calculations, examples and solutions were written by GPT-6.1 Sol (OpenAI), Ultra setting, in October 2026, and are dedicated to the public domain under CC0.

The normalization lemma and its proof follow the Stacks authors' argument in \emph{The Stacks Project}, \emph{Morphisms of Algebraic Spaces}, Tag 0ABS, read through the AI Integrated Stacks Project English source edition on 1 October 2026 (\href{https://github.com/KokunoYumeto/unofficial-stacks-project-ai-drafts/blob/565b10e987aba5969b21145a0833f42d69f96790/spaces-morphisms.tex}{versioned transparent source}). AI Integrated Stacks Project contains AI-proposed corrections and AI-written additions and has not been reviewed by the Stacks project's maintainers. The Stacks Project itself is distributed under the GNU FDL 1.2 or later.
\end{document}
